% From Hilbert Spaces to Sigma-Algebras: The Transformer Explained at a Higher Level of Mathematical Abstraction
% Included as the final chapter from main.tex; compiles with pdflatex.
\chapter{\texorpdfstring{From Hilbert Spaces to \(\sigma\)-Algebras: The Transformer Explained at a Higher Level of Mathematical Abstraction}{From Hilbert Spaces to Sigma-Algebras: The Transformer Explained at a Higher Level of Mathematical Abstraction}}
\chaptermark{From Hilbert Spaces to \(\sigma\)-Algebras}
\label{ch:math-transformers}

\section{Introduction}

A transformer is the composition of two maps of different mathematical
kinds. The first, from the input tokens to the last hidden state, is an
operator between finite-dimensional Hilbert spaces: nonlinear as a
whole, but assembled from linear maps, row-wise normalizations and
softmax-weighted averages. The second, the output layer, is a Markov
kernel from the last hidden state to the vocabulary: the softmax assigns
to each hidden state a probability measure on the \(\sigma\)-algebra of token
events, and the next token is a random variable with that law. This
chapter develops both maps for the original Transformer of 2017.

The chapter has three main sections.

\begin{itemize}
\tightlist
\item
  \textbf{Mathematical Tools} (Section~\ref{sec:math-tools}) collects the tools: Hilbert
  and Banach spaces, linear and nonlinear operators, spectral theory,
  weak derivatives and Sobolev spaces, measure-theoretic probability
  built on the \(\sigma\)-algebra of events, the softmax as a Gibbs measure,
  Fourier analysis, finite-precision arithmetic, information theory, and
  approximation theory.
\item
  \textbf{The Original Transformer} (Section~\ref{sec:math-original}) describes the 2017
  architecture as its authors presented it, with light notation.
\item
  \textbf{The Transformer as Operators and Measures}
  (Section~\ref{sec:math-operators}) writes the 2017 model as a composition of operators, then shows
  how its last layer defines a probability measure on the \(\sigma\)-algebra of
  token events and how tokens are sampled from it.
\end{itemize}

The chapter concerns what a transformer is: a parametrized family of
measures \(\{P_\theta\}\). How one member of the family is chosen, through
pretraining, post-training, alignment and reinforcement learning, is a
separate subject that this chapter does not treat.

\textbf{The map from architecture to mathematics.} Each component of the
model corresponds to one object of Section~\ref{sec:math-tools}. Table~\ref{tab:math-map} lists, for each
component of the architecture, the mathematical object and the
subsection that develops it. Every subsection of Section~\ref{sec:math-tools} appears on the right.

\begin{longtable}{@{}p{0.30\linewidth}p{0.47\linewidth}p{0.15\linewidth}@{}}
\caption{Each component of the 2017 Transformer, the mathematical object it corresponds to, and the subsection that develops it.}\label{tab:math-map}\\
\toprule
Component & Mathematical object & Developed in \\
\midrule
\endhead
Hidden states \(X\in\mathbb{R}^{n\times d}\) & Finite-dimensional Hilbert space, Frobenius inner product & \S\ref{sec:math:I1} \\
Layer stacking with \(\circ\) & Composition of operators, Lipschitz constants, chain rule & \S\ref{sec:math:I2} \\
Token embedding & Linear map into the Hilbert space & \S\ref{sec:math:I1} \\
Positional encoding & Orthogonal rotations, Fourier shift theorem & \S\ref{sec:math:I3}, \S\ref{sec:math:I8} \\
Attention scores \(QK^{\top}\) & Bilinear form, SVD, rank \(\le d_k\), Eckart--Young & \S\ref{sec:math:I3} \\
Attention softmax & Row-stochastic matrix, Markov kernel over positions & \S\ref{sec:math:I3}, \S\ref{sec:math:I6} \\
Cross-attention over the memory & Integral operator with the exponential kernel, compact operators, Mercer (in projected coordinates) & \S\ref{sec:math:I4} \\
Residual + LayerNorm & \(\mathrm{Id}+G\); projection theorem; compact set of states & \S\ref{sec:math:I1}, \S\ref{sec:math:I2} \\
Feed-forward ReLU & Piecewise-linear map, weak derivative, \(W^{1,\infty}\) & \S\ref{sec:math:I5} \\
Output logits & Linear functionals, Riesz representers & \S\ref{sec:math:I1} \\
Output softmax & Probability measure on the \(\sigma\)-algebra \(2^{\mathcal{V}}\), Gibbs measure & \S\ref{sec:math:I6}, \S\ref{sec:math:I7} \\
Temperature, top-\(k\), nucleus & Gibbs family, conditioning on an event & \S\ref{sec:math:I7} \\
Sampling a token & Random variable with that law, Gumbel-max & \S\ref{sec:math:I6}, \S\ref{sec:math:I7} \\
Causal mask & Predictability along a filtration, Doob--Dynkin & \S\ref{sec:math:I6} \\
Autoregressive loop & Measure on sequences, Ionescu-Tulcea, \(\pi\)-\(\lambda\) uniqueness & \S\ref{sec:math:I6} \\
End-of-sequence token & Stopping time, absorbing Markov chain & \S\ref{sec:math:I6} \\
Training target, label smoothing & Conditional expectation as projection, cross-entropy and KL & \S\ref{sec:math:I6}, \S\ref{sec:math:I10} \\
Floating-point evaluation & Finite precision, underflow of the measure & \S\ref{sec:math:I9} \\
Expressivity of the architecture & Stone--Weierstrass, universality & \S\ref{sec:math:I11} \\
\bottomrule
\end{longtable}

The same correspondence, read in the other direction, organizes Section~\ref{sec:math-operators}.

\textbf{Notation.} Vectors are rows when they are rows of a matrix
\(X \in \mathbb{R}^{n\times d}\), one row per token. In bilinear forms such as
\(x^{\top} B\, y\), and in matrix--vector products such as \(Wh\), the vectors involved
are read as columns.

\section{Mathematical Tools}\label{sec:math-tools}

This section builds every tool that Section~\ref{sec:math-operators} uses, in the order in which they
depend on each other. Finite-dimensional linear algebra comes first,
then its infinite-dimensional extensions, then measure-theoretic
probability, then the analytic tools for positions, precision,
information and approximation.

\subsection{Vector spaces, inner products, and Hilbert
spaces}\label{sec:math:I1}

A real vector space \(V\) is a set closed under addition and scalar
multiplication. An inner product \(\langle\cdot,\cdot\rangle\) on \(V\) equips it with geometry: it
induces a norm \(\|x\| = \sqrt{\langle x,x\rangle}\), a distance \(\|x - y\|\), and a notion of
orthogonality, \(\langle x,y\rangle = 0\). A Hilbert space is an inner-product space that
is complete: every Cauchy sequence converges to a point of the space.

Completeness is a nontrivial condition in infinite dimensions but holds
automatically in finite dimensions. Every finite-dimensional
inner-product space is a Hilbert space, isometrically isomorphic to
\(\mathbb{R}^{d}\) with the Euclidean inner product. The space of
matrices \(\mathbb{R}^{n\times d}\) with the Frobenius inner product is
a Hilbert space of dimension \(nd\):

\[
\langle X, Y\rangle_F = \operatorname{tr}(X^\top Y) = \sum_{i,j} X_{ij} Y_{ij}.
\]

In finite dimensions all norms are equivalent: for any two norms there
are constants \(c, C > 0\) with \(c\|x\|_a \le \|x\|_b \le C\|x\|_a\).
Convergence and continuity therefore do not depend on the norm, although
the constants can depend strongly on the dimension.

\textbf{The projection theorem.} Let \(M\) be a closed subspace of a Hilbert
space \(H\) (in finite dimensions every subspace is closed). Every \(x \in H\)
decomposes uniquely as

\[
x = P_M x + (x - P_M x), \qquad P_M x \in M, \quad x - P_M x \perp M,
\]

and \(P_M x\) is the unique point of \(M\) closest to \(x\). The map \(P_M\) is
linear, idempotent and self-adjoint. Least squares, and in \S\ref{sec:math:I6}
conditional expectation, are instances of this theorem.

\textbf{The Riesz representation theorem.} Every continuous linear
functional \(f\) on a Hilbert space \(H\) has the form \(f(x) = \langle x, y_f\rangle\) for a
unique \(y_f \in H\), and \(\|f\| = \|y_f\|\). In \(\mathbb{R}^{d}\) this says that
every linear functional is \(x \mapsto u \cdot x\) for a unique vector \(u\), its
representer. A vector can therefore be read in two ways: as a point of
the space, or as the functional it represents.

\subsection{Operators: linear, nonlinear, and
piecewise}\label{sec:math:I2}

An operator is a map \(T\colon V \to W\) between vector spaces. It is linear if
\(T(\alpha x + \beta y) = \alpha T(x) + \beta T(y)\). Between finite-dimensional spaces, every
linear operator is represented by a matrix once bases are fixed, and is
bounded, with operator norm

\[
\|T\| = \sup_{\|x\| = 1} \|Tx\|,
\]

whose value depends on the norms chosen on \(V\) and \(W\).

Nonlinear operators are studied through weaker structures.

\begin{itemize}
\tightlist
\item
  \textbf{Lipschitz continuity.} \(T\) is \(L\)-Lipschitz if \(\|T(x) - T(y)\| \le L\|x - y\|\)
  for all \(x, y\). It is locally Lipschitz if such a bound holds on
  every bounded set, with a constant that may grow with the set.
\item
  \textbf{Differentiability.} \(T\) is Fréchet differentiable at \(x\) if there
  is a linear operator \(DT(x)\), the Jacobian in coordinates, such that
  \(T(x + h) = T(x) + DT(x)h + o(\|h\|)\). Differentiation replaces a nonlinear
  operator locally by a linear one.
\end{itemize}

Both properties compose well: Lipschitz constants multiply, \(L(T\circ S) \le
L(T)\cdot L(S)\), and derivatives obey the chain rule, \(D(T\circ S)(x) = DT(S(x))\,
DS(x)\). Backpropagation is the chain rule applied, in reverse order, to a
long composition.

\textbf{Piecewise-smooth maps.} A map \(T\colon \mathbb{R}^{m} \to
\mathbb{R}^{p}\) is piecewise smooth if there are finitely many
disjoint open sets \(\Omega_1, \ldots, \Omega_r\) whose closures cover
\(\mathbb{R}^{m}\), \(T\) is smooth on each \(\Omega_j\), and the boundary set \(\Gamma\) has
Lebesgue measure zero. Such a map is differentiable almost everywhere.
It need not be continuous across \(\Gamma\), and when continuous it need not be
Lipschitz.

Two examples delimit the class. The ReLU function \(x \mapsto \max(0, x)\) is
piecewise linear, continuous and 1-Lipschitz. The top-\(k\) selection
operator, which returns the index set \(S_k(s)\) of the \(k\) largest entries
of a score vector \(s \in \mathbb{R}^{N}\), is piecewise constant and
jumps where two scores tie. A map of the form

\[
F(x) = \sum_{i \in S_k(s(x))} g_i(x)
\]

is smooth away from the tie set \(\{x : s_i(x) = s_j(x) \text{ for some } i \neq j\}\)
and generally discontinuous on it. The tie set has measure zero whenever
the score differences \(s_i - s_j\) are non-constant real-analytic
functions, because the zero set of a non-constant real-analytic function
has Lebesgue measure zero. This is the real-variable counterpart of the
identity theorem for holomorphic functions, and it applies to networks
built from analytic activations such as exp and GELU.

Rademacher's theorem states that every locally Lipschitz map is
differentiable almost everywhere. It covers ReLU networks but not
networks containing hard selections.

\textbf{Absolute continuity and the fundamental theorem of calculus.}
Almost-everywhere differentiability does not recover a function from its
derivative. On an interval, the identity \(f(b) - f(a) = \int_a^b f'\)
holds exactly when \(f\) is absolutely continuous, and every Lipschitz
function is absolutely continuous. The Cantor--Vitali function is the
standard counterexample: it is continuous and non-decreasing from 0 to
1, yet \(f' = 0\) almost everywhere. Its almost-everywhere derivative
carries no information about how the function changes. Step functions
fail more severely: they are discontinuous, and their derivative is zero
wherever it exists. For quantizers and hard selection, gradients computed almost
everywhere are therefore uninformative, and training requires surrogate
gradients.

\subsection{Spectral theory in finite
dimensions}\label{sec:math:I3}

\textbf{The spectral theorem.} A linear operator \(A\) on a real
inner-product space is self-adjoint if \(\langle Ax, y\rangle = \langle x, Ay\rangle\). Every
self-adjoint operator on a finite-dimensional space has an orthonormal
basis of eigenvectors with real eigenvalues:

\[
A = Q \Lambda Q^\top, \qquad Q^\top Q = I, \quad \Lambda = \operatorname{diag}(\lambda_1, \dots, \lambda_d).
\]

The theorem requires linearity and symmetry and says nothing about
nonlinear maps.

\textbf{Orthogonal operators.} A matrix \(Q\) with \(Q^{\top}Q = I\) is an
isometry, \(\|Qx\| = \|x\|\), and all its singular values equal 1. A
block-diagonal matrix of \(2 \times 2\) rotations \(R_\theta\) is the basic example.
Rotations satisfy \(R_\alpha R_\beta = R_{\alpha+\beta}\): composing them adds angles, so
they represent translation linearly. Identifying \(\mathbb{R}^{2}\) with
\(\mathbb{C}\), \(R_\theta\) is multiplication by \(e^{i\theta}\) (\S\ref{sec:math:I8}).

\textbf{The singular value decomposition.} Any \(W \in
\mathbb{R}^{m\times n}\) of rank \(r\) factors as

\[
W = U \Sigma V^\top = \sum_{i=1}^{r} \sigma_i\, u_i v_i^\top, \qquad \sigma_1 \ge \cdots \ge \sigma_r > 0,
\]

with orthonormal singular vectors. The SVD can be derived from the
spectral theorem applied to \(W^{\top}W\), whose nonzero eigenvalues are
the \(\sigma_i^{2}\). The largest singular value is the Euclidean operator
norm, \(\|W\|_{2} = \sigma_{1}\), and hence the Euclidean Lipschitz
constant of \(x \mapsto Wx\).

\textbf{Low-rank approximation.} By the Eckart--Young--Mirsky theorem,
truncating the SVD after \(k\) terms gives the best rank-\(k\) approximation in
both the spectral and Frobenius norms:

\[
\|W - W_k\|_2 = \sigma_{k+1}, \qquad \|W - W_k\|_F = \Big(\sum_{i > k} \sigma_i^2\Big)^{1/2}.
\]

Any product \(W = BA\) with \(B \in \mathbb{R}^{m\times k}\) and \(A \in
\mathbb{R}^{k\times n}\) has rank at most \(k\), so factoring through a
\(k\)-dimensional space imposes a rank bound by construction.

\textbf{The polar decomposition.} For \(W \in \mathbb{R}^{m\times n}\)
with \(m \ge n\), the thin SVD gives

\[
W = (U V^\top)\,(V \Sigma V^\top),
\]

the product of a matrix with orthonormal columns and a positive
semidefinite matrix. The polar factor \(UV^{\top}\) is the nearest
matrix with orthonormal columns in the Frobenius norm: it keeps the
directions of \(W\) and sets every singular value to 1. Newton--Schulz
iterations compute it using matrix products only.

\textbf{Stochastic matrices.} \(P \in \mathbb{R}^{n\times m}\) is
row-stochastic if \(P \ge 0\) entrywise and every row sums to 1. For any \(V \in
\mathbb{R}^{m\times d}\), each row of \(PV\) is a convex combination of
rows of \(V\), and Hölder's inequality with exponents \(1\) and \(\infty\) gives, for
every column \(j\),

\[
|(PV)_{ij}| \le \sum_k P_{ik}\,|V_{kj}| \le \max_k |V_{kj}|.
\]

A row-stochastic matrix is therefore non-expansive in the maximum norm,
\(\|P\|_{\infty\to\infty} = 1\). It need not be non-expansive in the Euclidean norm: the
\(n \times n\) matrix whose rows all equal \(e_{1}^{\top}\) has \(\|P\|_{2} = \sqrt{n}\).

\textbf{Doubly stochastic matrices.} A square matrix \(B\) is doubly
stochastic if both \(B\) and \(B^{\top}\) are row-stochastic. By the
Birkhoff--von Neumann theorem, the doubly stochastic matrices form the
convex hull of the permutation matrices (the Birkhoff polytope). Two
consequences follow:

\begin{itemize}
\tightlist
\item
  every doubly stochastic matrix satisfies \(\|B\|_{2} \le 1\), since each
  permutation matrix has norm 1 and the norm is convex;
\item
  the polytope is closed under multiplication, so products of
  arbitrarily many doubly stochastic matrices remain non-expansive.
\end{itemize}

By the Sinkhorn--Knopp theorem, for every matrix \(A\) with strictly
positive entries there are positive diagonal matrices \(D_{1}\),
\(D_{2}\), unique up to a scalar, such that \(D_{1}AD_{2}\) is
doubly stochastic. Alternating row and column normalization converges to
it; stopping after finitely many steps gives an approximately doubly
stochastic matrix.

\subsection{Banach spaces, compact operators, and the spectral
theorem}\label{sec:math:I4}

Two later topics need infinite dimensions. Approximation theory (\S\ref{sec:math:I11})
works in spaces of functions that are generally not Hilbert spaces, and
the kernel view of attention (Section~\ref{sec:math-operators}) describes attention as an
integral operator, whose structure is governed by the spectral theory of
compact operators.

\textbf{Banach spaces.} A normed space is a Banach space if it is
complete. Every Hilbert space is Banach; a norm comes from an inner
product exactly when it satisfies the parallelogram law
\(\|x + y\|^{2} + \|x - y\|^{2} = 2\|x\|^{2} + 2\|y\|^{2}\). Two examples recur:

\begin{itemize}
\tightlist
\item
  \(C(K)\), the continuous functions on a compact \(K \subset \mathbb{R}^{d}\)
  with the sup norm \(\|f\|_\infty = \max_{K} |f|\): Banach but
  not Hilbert, and the natural setting for uniform approximation;
\item
  \(L^p(\mu)\) for \(1 \le p \le \infty\), with
  \(\|f\|_p = \bigl(\int |f|^p \, d\mu\bigr)^{1/p}\): Hilbert only for \(p = 2\).
  Hölder's inequality \(\|fg\|_{1} \le \|f\|_p\|g\|_q\) (\(1/p + 1/q = 1\))
  and Minkowski's inequality \(\|f + g\|_p \le \|f\|_p + \|g\|_p\) are its basic tools.
\end{itemize}

Two finite-dimensional facts fail in infinite dimensions.

\begin{enumerate}
\def\labelenumi{\arabic{enumi}.}
\tightlist
\item
  \textbf{Norms are no longer equivalent.} On \(C([0,1])\) there are
  sequences \((f_n)\) with \(\|f_n\|_{1}\) constant and \(\|f_n\|_\infty \to \infty\). A claim
  such as ``networks approximate \(f\) to accuracy \(\varepsilon\)'' is incomplete until
  the norm is named.
\item
  \textbf{Bounded sets need not have compact closure.} By Riesz's lemma,
  the closed unit ball is compact if and only if the space is
  finite-dimensional. Compactness, automatic in \(\mathbb{R}^{d}\) by
  Heine--Borel, becomes a special property.
\end{enumerate}

\textbf{Bounded and compact operators.} A linear operator \(T\colon X \to Y\) is
bounded if \(\|T\| = \sup_{\|x\| \le 1} \|Tx\|\) is finite, which is equivalent
to continuity; when \(Y\) is Banach, the bounded operators \(B(X, Y)\) form a
Banach space. \(T\) is compact if it maps bounded sets to sets with compact
closure. Finite-rank operators are compact, and on Hilbert spaces the
compact operators are exactly the operator-norm limits of finite-rank
operators.

The main example is the integral operator. For a finite Borel measure \(\mu\)
on a compact \(K\) and a kernel \(\kappa \in L^{2}(\mu \otimes \mu)\),

\[
(T_\kappa f)(u) = \int_K \kappa(u, v)\, f(v)\, d\mu(v)
\]

is a Hilbert--Schmidt, hence compact, operator on \(L^{2}(\mu)\). Every
continuous kernel on a compact set defines one.

\textbf{The spectral theorem for compact self-adjoint operators.} Let \(T\)
be compact and self-adjoint on a separable Hilbert space \(H\). There are an
orthonormal sequence \((e_i)\) and real eigenvalues \(\lambda_i\) with
\(|\lambda_{1}| \ge |\lambda_{2}| \ge \cdots\) and \(\lambda_i
\to 0\) such that

\[
Tf = \sum_i \lambda_i\, \langle f, e_i \rangle\, e_i \qquad \text{for all } f \in H,
\]

and each nonzero eigenvalue has a finite-dimensional eigenspace. This is
the infinite-dimensional analogue of \(A = Q\Lambda Q^{\top}\), with the
additional constraint that the eigenvalues can accumulate only at zero.
For compact operators that are not self-adjoint, the Schmidt
decomposition \(Tf = \sum s_i \langle f, v_i\rangle u_i\) plays the role of the SVD, with
singular values \(s_i\) decreasing to zero. Eckart--Young carries over:
truncating after \(k\) terms is optimal, with error \(s_{k+1}\) in operator
norm. The decay of the singular values measures how compressible the
operator is.

\textbf{Mercer's theorem.} A symmetric kernel \(\kappa\) is positive definite if
every matrix \([\kappa(u_i, u_j)]\) built from finitely many points is
positive semidefinite. Let \(K\) be compact, \(\mu\) a finite Borel measure with
support \(K\), and \(\kappa\) continuous, symmetric and positive definite. Then \(T_\kappa\)
is compact, self-adjoint and positive, and

\[
\kappa(u, v) = \sum_i \lambda_i\, \phi_i(u)\, \phi_i(v), \qquad \lambda_i \ge 0,
\]

with absolute and uniform convergence on \(K \times K\). The kernel is an inner
product of feature vectors \(\Phi(u) = (\sqrt{\lambda_i}\, \phi_i(u))_i\) in \(\ell^{2}\),
and these features span its reproducing kernel Hilbert space (RKHS).

The exponential kernel \(\kappa(q, k) = e^{\langle q,k\rangle}\) is positive definite: its
Taylor series is a sum, with nonnegative coefficients, of powers of the
linear kernel, and powers of a positive-definite kernel remain positive
definite by the Schur product theorem. On a compact set Mercer's theorem
applies to it. Because the kernel is analytic, its eigenvalues decay
rapidly; the decay slows as the domain grows, because the kernel then
varies over a wider range of scales. The precise rate depends on the dimension and the domain
radius; see Steinwart and Christmann for quantitative statements.

\subsection{Weak derivatives and Sobolev
spaces}\label{sec:math:I5}

The states of a transformer are vectors in finite-dimensional spaces, so
Sobolev spaces do not describe the state space. They are the right
language for two other questions: how regular the network map itself is,
and how regular a target function must be for networks to approximate it
at a given rate (\S\ref{sec:math:I11}).

\textbf{Weak derivatives and distributions.} Let \(\Omega \subset \mathbb{R}^{m}\)
be open and \(f\) locally integrable. A locally integrable \(g\) is the weak
partial derivative \(\partial_i f\) if

\[
\int_\Omega f\, \partial_i \varphi \, dx = -\int_\Omega g\, \varphi \, dx \qquad \text{for all } \varphi \in C_c^\infty(\Omega).
\]

More generally, a distribution is a continuous linear functional \(T\) on
the test functions \(C_c^\infty(\Omega)\), and its derivative is defined by
\(\langle \partial_i T, \varphi\rangle = -\langle T, \partial_i \varphi\rangle\). Every locally integrable function defines a
regular distribution \(\varphi \mapsto \int f\varphi\), and every distribution has derivatives of
all orders.

The model example is ReLU. Its weak derivative is the Heaviside step \(H =
\mathbf{1}_{(0,\infty)}\), and its second distributional derivative is the Dirac delta
\(\delta_{0}\), the functional \(\varphi \mapsto \varphi(0)\). A network built from ReLU and linear
maps is piecewise linear, so its classical Hessian is zero almost
everywhere; its curvature is concentrated, as a measure, on the
piecewise-linear hypersurfaces (hyperplanes for a single layer, bent
hyperplanes for deeper networks) where some unit switches.

\textbf{Sobolev spaces.} For \(k \in \mathbb{N}\) and \(1 \le p \le \infty\),

\[
W^{k,p}(\Omega) = \{ f \in L^p(\Omega) : \partial^\alpha f \in L^p(\Omega) \ \text{for all } |\alpha| \le k \},
\]

with weak derivatives, is a Banach space; \(H^k = W^{k,2}\) is a
Hilbert space. Two facts connect this to \S\ref{sec:math:I2}.

\begin{itemize}
\tightlist
\item
  \textbf{Lipschitz functions are \(W^{1,\infty}\).} On \(\mathbb{R}^{m}\), a
  function is locally Lipschitz exactly when it belongs to
  \(W^{1,\infty}_{\mathrm{loc}}\), and its weak gradient coincides with the classical
  gradient that Rademacher's theorem provides almost everywhere. Dense
  attention, smooth activations, ReLU and normalization (smooth everywhere for \(\epsilon > 0\)) all produce network maps in this class.
\item
  \textbf{Jumps leave the Sobolev scale.} A function that jumps across a
  hypersurface is not in \(W^{1,1}_{\mathrm{loc}}\). It is locally of bounded
  variation: its distributional derivative is a measure with an
  absolutely continuous part (the almost-everywhere gradient) and a
  singular jump part concentrated on the hypersurface. Hard top-\(k\)
  selection, discrete routing and quantizers produce maps of this kind.
  The almost-everywhere gradient sees only the first part and misses the
  jumps entirely, which is the precise form of the warning at the end of
  \S\ref{sec:math:I2}.
\end{itemize}

\subsection{\texorpdfstring{Measure, probability, and the \(\sigma\)-algebra of
events}{Measure, probability, and the sigma-algebra of events}}\label{sec:math:I6}

This section builds the probabilistic language of the chapter. Its
central object is the \(\sigma\)-algebra: the collection of events to which
probabilities can be consistently assigned.

\textbf{\(\sigma\)-algebras and events.} A \(\sigma\)-algebra \(\mathcal{F}\) on a set \(\Omega\) is a collection
of subsets that contains \(\Omega\) and is closed under complements and countable
unions. The pair \((\Omega, \mathcal{F})\) is a measurable space, and its elements are
events. The axioms guarantee that ``not'', ``or'' and ``and'', applied
to countably many events, again yield events. For any collection \(\mathcal{C}\) of
subsets, \(\sigma(\mathcal{C})\) denotes the smallest \(\sigma\)-algebra containing \(\mathcal{C}\), the \(\sigma\)-algebra
it generates.

\textbf{The token \(\sigma\)-algebra.} When \(\Omega\) is finite or countable, the power
set \(2^{\Omega}\) is the natural \(\sigma\)-algebra. For a finite vocabulary \(\mathcal{V}\), the
space \((\mathcal{V}, 2^{\mathcal{V}})\) has \(2^{|\mathcal{V}|}\) events; its atoms are
the singletons \(\{v\}\), and every event is a finite union of atoms. Events
have direct linguistic readings: ``the next token is punctuation'',
``the next token begins a word'', ``the next token lies in a given top-\(k\)
set''.

\textbf{Borel \(\sigma\)-algebras.} On \(\mathbb{R}^{d}\) the power set is not a
usable \(\sigma\)-algebra, since Lebesgue measure cannot be extended to all
subsets. The standard choice is the Borel \(\sigma\)-algebra
\(\mathcal{B}(\mathbb{R}^{d})\), generated by the open sets. Continuous and
piecewise-smooth functions are Borel measurable. Hidden states, and
continuous inputs such as image embeddings, live in measurable spaces of
this kind.

\textbf{Measures, conditioning, and construction.} A probability measure
\(P\) on \((\Omega, \mathcal{F})\) is countably additive with \(P(\Omega) = 1\). On a countable space it
is determined by its values on atoms, \(P(S) = \sum_{\omega \in S} p(\omega)\). For
an event \(S\) with \(P(S) > 0\), conditioning gives the measure
\(P(A \mid S) = P(A \cap S)/P(S)\), supported on \(S\). A sub-probability
measure satisfies \(P(\Omega) \le 1\); adjoining a cemetery point \(\partial\) that carries
the missing mass turns it into a probability measure.

Measures on large \(\sigma\)-algebras are built from their values on simple
generating classes. Carathéodory's construction extends a countably
additive set function from an algebra to the \(\sigma\)-algebra it generates.
Uniqueness comes from Dynkin's \(\pi\)-\(\lambda\) theorem: a \(\pi\)-system is a class closed
under finite intersections, and two probability measures that agree on a
\(\pi\)-system agree on the \(\sigma\)-algebra it generates.

\textbf{Random variables, information, and Doob--Dynkin.} A map \(X\colon (\Omega,
\mathcal{F}) \to (E, \mathcal{E})\) is measurable if \(X^{-1}(B) \in \mathcal{F}\) for every \(B \in \mathcal{E}\); on a
probability space it is a random variable. Its law \(P^X(B) = P(X \in B)\)
is the pushforward of \(P\). Each random variable generates \(\sigma(X) =
\{X^{-1}(B) : B \in \mathcal{E}\}\), the events decided by the value of \(X\): its
information. The Doob--Dynkin lemma makes this precise. If \(Y\) takes
values in \(\mathbb{R}^{d}\) or a finite set, then \(Y\) is \(\sigma(X)\)-measurable
if and only if \(Y = f(X)\) for some measurable \(f\).~``\(Y\) depends only on \(X\)''
and ``\(Y\) is measurable with respect to the information in \(X\)'' are the
same statement.

\textbf{Densities and Radon--Nikodym.} A measure \(\nu\) is absolutely
continuous with respect to \(\mu\), written \(\nu \ll \mu\), if \(\mu(A) = 0\) implies \(\nu(A) =
0\). For \(\sigma\)-finite measures, the Radon--Nikodym theorem gives a density
\(d\nu/d\mu \ge 0\) with \(\nu(A) = \int_A (d\nu/d\mu)\, d\mu\) for every \(A\). On a finite space the
density is simply the ratio of masses, \(\nu(\{v\})/\mu(\{v\})\). Likelihood
ratios and the Kullback--Leibler divergence (\S\ref{sec:math:I10}) are defined through
Radon--Nikodym derivatives.

\textbf{Conditional expectation as a projection.} Let \(\mathcal{G} \subset \mathcal{F}\) be a
sub-\(\sigma\)-algebra. For \(X \in L^{1}(P)\), the conditional expectation \(\mathbb{E}[X
\mid \mathcal{G}]\) is the almost surely unique \(\mathcal{G}\)-measurable random variable
with

\[
\int_A \mathbb{E}[X \mid \mathcal{G}]\, dP = \int_A X \, dP \qquad \text{for all } A \in \mathcal{G}.
\]

For \(X \in L^{2}\), the \(\mathcal{G}\)-measurable square-integrable variables form a
closed subspace \(L^{2}(\Omega, \mathcal{G}, P)\) of the Hilbert space \(L^{2}(\Omega, \mathcal{F},
P)\), and \(\mathbb{E}[X \mid \mathcal{G}]\) is exactly the orthogonal projection of \(X\)
onto it (\S\ref{sec:math:I1}). It is therefore the best predictor of \(X\) in mean square
among all functions of the information in \(\mathcal{G}\). Conditional probabilities
are the case \(X = 1_A\): \(P(A \mid \mathcal{G}) = \mathbb{E}[1_A \mid \mathcal{G}]\).

\textbf{Kernels.} A Markov kernel from \((E, \mathcal{E})\) to \((\Omega, \mathcal{F})\) assigns to each
\(e \in E\) a probability measure \(K(e, \cdot)\) on \(\mathcal{F}\), measurably in \(e\). It is a
measure that varies with an input, the formal object behind every
conditional distribution in this chapter. A row-stochastic matrix is
exactly a Markov kernel between finite sets; if \(K(e, \Omega) \le 1\) is allowed,
the kernel is sub-Markov.

\textbf{Sequence space, cylinder sets, and filtrations.} Let \(\Omega =
\mathcal{V}^{\mathbb{N}}\) be the space of infinite sequences over a finite
alphabet, with product \(\sigma\)-algebra. The cylinder sets

\[
C(x_1, \dots, x_t) = \{ \omega \in \Omega : \omega_1 = x_1, \dots, \omega_t = x_t \}
\]

fix a prefix. For each \(t\), \(\mathcal{F}_{t} = \sigma(\omega_{1}, \ldots, \omega_{t})\) is
generated by the cylinders of length \(t\) and contains every event decided
by the first \(t\) symbols. The \(\sigma\)-algebras are nested, \(\mathcal{F}_{0} = \{\emptyset, \Omega\}
\subset \mathcal{F}_{1} \subset \mathcal{F}_{2} \subset \cdots\), and form a filtration; \(\mathcal{F}_\infty = \sigma(\bigcup_{t}
\mathcal{F}_{t})\) is the product \(\sigma\)-algebra. The cylinders together with \(\emptyset\) form
a \(\pi\)-system generating \(\mathcal{F}_\infty\), so by the \(\pi\)-\(\lambda\) theorem a probability measure
on sequence space is uniquely determined by the probabilities of all
prefixes.

\textbf{Adaptedness, predictability, and stopping times.} A process
\((Y_t)\) is adapted if each \(Y_t\) is \(\mathcal{F}_t\)-measurable, and predictable if
each \(Y_t\) is \(\mathcal{F}_{t-1}\)-measurable: determined strictly by the past. By
Doob--Dynkin, predictability means \(Y_t = f_t(\omega_{1}, \ldots,
\omega_{t-1})\). A stopping time is a random time \(\tau\) with \(\{\tau = t\} \in \mathcal{F}_t\) for
every \(t\): whether to stop is decidable from the information available at
that time.

\textbf{Construction of sequence measures.} The Ionescu-Tulcea theorem
states that a sequence of Markov kernels \(p(\cdot \mid
x_{<t})\) from the past to the next symbol determines a unique
probability measure on \((\Omega, \mathcal{F}_\infty)\). On cylinders it satisfies the chain
rule:

\[
P\big(C(x_1, \dots, x_t)\big) = \prod_{s=1}^{t} p(x_s \mid x_{<s}).
\]

When generation starts from a prompt, the construction applies
conditionally on it.

\textbf{Markov chains and absorption.} Such a process is a
time-homogeneous Markov chain on the countable space of finite
prefixes (a prefix records its own length, so no separate time index is
needed): each step appends one symbol according to a kernel that
depends on the current prefix. A state is absorbing if the chain, once
there, never leaves. Absorption problems ask for the probability that
the chain is eventually absorbed; Section~\ref{sec:math-operators} identifies this probability
with termination of generation.

\subsection{The softmax and Gibbs
measures}\label{sec:math:I7}

For \(z \in \mathbb{R}^{K}\) and a temperature \(\tau > 0\), the
softmax

\[
\operatorname{softmax}_\tau(z)_k = \frac{e^{z_k/\tau}}{\sum_j e^{z_j/\tau}}
\]

maps any real vector to a strictly positive probability vector, and
hence to a probability measure on \((\{1, \ldots, K\},
2^{\{1,\ldots,K\}})\) that charges every atom.

\begin{itemize}
\tightlist
\item
  \textbf{Shift invariance.} \(\operatorname{softmax}_\tau(z + c\mathbf{1}) =
  \operatorname{softmax}_\tau(z)\) for every constant \(c\).~Only logit differences matter; the
  logits live naturally in the quotient
  \(\mathbb{R}^{K}/\mathbb{R}\mathbf{1}\).
\item
  \textbf{Gradient and curvature.} The softmax is the gradient of the
  convex function \(\tau \log \sum_j e^{z_j/\tau}\). At \(\tau = 1\) its Jacobian, which
  is also the Hessian of log-sum-exp, is \(\operatorname{diag}(p) - pp^{\top}\). This
  is the covariance matrix of the one-hot vector drawn from \(p\), and the
  Fisher information of the categorical family in its natural
  parameters. It is symmetric positive semidefinite with \(\mathbf{1}\)
  in its null space, and \(\operatorname{softmax}_\tau\) is \((1/\tau)\)-Lipschitz in the Euclidean
  norm.
\item
  \textbf{Variational characterization.} The softmax is the Gibbs
  distribution:
\end{itemize}

\[
\operatorname{softmax}_\tau(z) = \arg\max_{p \in \Delta^{K-1}} \big\{ \langle p, z \rangle + \tau H(p) \big\}, \qquad H(p) = -\sum_k p_k \log p_k.
\]

This is convex duality: the Fenchel conjugate of log-sum-exp is negative
entropy on the simplex and \(+\infty\) off it. Temperature weights entropy
against expected score. As \(\tau \to 0\) the measure converges weakly to the
uniform measure on the argmax set, a Dirac mass when the maximizer is
unique; as \(\tau \to \infty\) it tends to the uniform measure.

\begin{itemize}
\tightlist
\item
  \textbf{Masking and restriction.} Logits set to \(-\infty\) receive probability
  zero. Applying the softmax only to the logits indexed by a set \(S\)
  produces exactly the conditional measure \(P(\cdot \mid S)\) of \S\ref{sec:math:I6}:
  restricting the softmax to a subset is conditioning on that event.
\item
  \textbf{Extra denominator mass.} Adding \(b > 0\) to the
  denominator yields a sub-probability vector, equivalently a
  probability vector on \(\{1, \ldots, K\} \cup \{\partial\}\) with mass \(b/(b + \sum_j
  e^{z_j})\) on the cemetery point.
\item
  \textbf{Sampling as a measurable map of noise.} If \(G_{1}, \ldots,
  G_K\) are independent standard Gumbel variables, then \(\arg\max_k (z_k +
  G_k)\) has law \(\operatorname{softmax}_{1}(z)\) (the Gumbel-max trick).
\item
  \textbf{Stable computation.} Shift invariance justifies computing the
  softmax after subtracting \(\max_k z_k\), which prevents overflow.
  Partial results over blocks combine associatively: a block summarized
  by \((m, s)\), with \(m\) its maximum logit and \(s\) the sum of \(e^{z_k - m}\),
  merges with another as \(m = \max(m_{1}, m_{2})\), \(s =
  s_{1}e^{m_{1}-m} + s_{2}e^{m_{2}-m}\). This online
  softmax lets attention be computed block by block without storing the
  full score matrix.
\end{itemize}

\subsection{Fourier analysis and positional
structure}\label{sec:math:I8}

\textbf{Complex exponentials and rotations.} Identify \(\mathbb{R}^{2}\)
with \(\mathbb{C}\). The rotation \(R_\theta\) becomes multiplication by
\(e^{i\theta}\), and the map \(m \mapsto e^{im\theta}\) turns translation of an integer
position \(m\) into multiplication: \(e^{i(m+k)\theta} =
e^{ik\theta}e^{im\theta}\). This identity underlies both Fourier
analysis and positional encodings.

\textbf{Fourier series.} The functions \(e_k(x) = e^{ikx}/\sqrt{2\pi}\), \(k \in
\mathbb{Z}\), form an orthonormal basis of \(L^{2}(-\pi, \pi)\). Every \(f \in
L^{2}\) expands as \(f = \sum \langle f, e_k\rangle e_k\), and Parseval's identity
\(\|f\|^{2} = \sum |\langle f, e_k\rangle|^{2}\) says the coefficient
map is an isometry onto \(\ell^{2}\).

\textbf{The Fourier transform.} For \(f \in L^{1}(\mathbb{R}^{d})\),
\(\hat{f}(\xi) = \int f(x)e^{-ix\cdot\xi}\, dx\). It satisfies the Riemann--Lebesgue lemma,
inversion for suitable \(f\), and, by Plancherel's theorem, extends to a
multiple of a unitary operator on \(L^{2}\). Two properties are used
below:

\[
\widehat{f(\cdot - a)}(\xi) = e^{-i a \cdot \xi}\, \hat f(\xi), \qquad \widehat{f * g} = \hat f\, \hat g.
\]

The shift theorem says translation acts on frequencies as a phase
rotation. The convolution theorem says every bounded translation-invariant
linear operator on \(L^2\), which is a convolution (in general with a
distribution), is diagonalized by the Fourier transform.

\textbf{Phases and their ambiguity.} The argument of a complex number is
defined only modulo \(2\pi\); the complex logarithm is multivalued, and a
branch must be chosen. For positional structure this means that the
phase \(m\theta\) determines the position \(m\) only modulo \(2\pi/\theta\). A single frequency
cannot distinguish positions that differ by a full period (aliasing).
Several frequencies in geometric progression resolve positions over a
range set by the longest period. Incommensurate frequencies never alias
exactly on the integers, but beyond that range distinct positions become
numerically hard to tell apart.

\subsection{Finite-precision
arithmetic}\label{sec:math:I9}

Digital hardware does not compute in \(\mathbb{R}\). A floating-point
format defines a finite set \(\mathbb{F} \subset \mathbb{R} \cup \{\pm\infty\}\),
and rounding is a map \(\mathrm{fl}\colon \mathbb{R} \to \mathbb{F}\) sending each
real number to a nearby element. Within the normal range \(|\mathrm{fl}(x)
- x| \le u|x|\), where the unit roundoff \(u\) is set by
the number of mantissa bits.

Low-precision formats combine a small grid with a shared scale: a block
of \(b\) values is stored as \(x_{i} \approx s \cdot q_{i}\), each \(q_{i}\) from
the grid and one scale \(s\) per block. The 4-bit E2M1 format has magnitudes
\(\{0, 0.5, 1, 1.5, 2, 3, 4, 6\}\) and a sign bit, so the largest value in a
block is \(6 \cdot s_{\max}\), where \(s_{\max}\) is the largest representable scale.

A quantizer \(Q = \mathrm{fl}\) has three properties:

\begin{itemize}
\tightlist
\item
  it is idempotent, \(Q \circ Q = Q\);
\item
  it is piecewise constant on each coordinate, hence nonlinear and
  locally of bounded variation;
\item
  its derivative vanishes except on the finite set of rounding
  thresholds, where it jumps.
\end{itemize}

By \S\ref{sec:math:I2} and \S\ref{sec:math:I5}, its almost-everywhere derivative carries no
information, so training through a quantizer needs a surrogate gradient,
commonly the straight-through estimator, which treats \(DQ\) as the
identity.

Norm bounds determine whether a format can overflow. If normalization
bounds the Euclidean norm of a state by \(R\), every coordinate is bounded
by \(R\), and orthogonal maps such as rotations preserve that bound. A
format whose largest magnitude exceeds \(R\) cannot overflow on those
states. The normalized states lie on a compact set, so continuous maps
of them are bounded (Weierstrass) and uniformly continuous
(Heine--Cantor).

The same remarks apply to probabilities. The exact softmax is strictly
positive, but in floating point small probabilities can underflow to
zero: the computed measure can assign zero mass to atoms that the
idealized measure charges.

\subsection{Information theory}\label{sec:math:I10}

Let \(p\) and \(q\) be probability vectors on a finite set. The entropy,
cross-entropy and Kullback--Leibler divergence are

\[
H(p) = -\sum_k p_k \log p_k, \qquad H(p, q) = -\sum_k p_k \log q_k, \qquad \mathrm{KL}(p \,\|\, q) = \sum_k p_k \log \frac{p_k}{q_k}.
\]

They satisfy \(H(p, q) = H(p) + \mathrm{KL}(p \,\|\, q)\). The divergence is the expected
log Radon--Nikodym derivative, \(\mathrm{KL}(P \,\|\, Q) = \mathbb{E}_P[\log dP/dQ]\), defined
when \(P \ll Q\) and \(+\infty\) otherwise. By Gibbs' inequality \(\mathrm{KL}(p \,\|\, q) \ge 0\), with
equality exactly when \(p = q\); minimizing cross-entropy in \(q\) therefore
recovers \(p\).~KL is not symmetric: when \(q\) ranges over a restricted family (for
example unimodal \(q\) and multimodal \(p\)), minimizing \(\mathrm{KL}(p \,\|\, q)\) over \(q\) spreads \(q\)
over all of \(p\)'s mass, while minimizing \(\mathrm{KL}(q \,\|\, p)\) concentrates \(q\) on a mode
of \(p\).

For measures on sequence space built by the chain rule, divergences
decompose along the filtration:

\[
\mathrm{KL}(P \,\|\, Q) = \sum_{t \ge 1} \mathbb{E}_P\Big[\mathrm{KL}\big(P(\cdot \mid x_{<t}) \,\big\|\, Q(\cdot \mid x_{<t})\big)\Big].
\]

The perplexity of a model \(q\) on a text \(x_{1}, \ldots, x_T\) is
\(\exp\bigl(-(1/T) \sum_{t} \log q(x_{t} \mid \text{earlier tokens})\bigr)\): the
exponential of the average per-token cross-entropy. A perplexity of \(k\)
means the model is, on average, as uncertain as a uniform choice among \(k\)
tokens.

\subsection{Approximation theory}\label{sec:math:I11}

Approximation theory asks which functions a class of networks can
represent, and at what cost. Its results are existence statements: they
say nothing about whether training finds the approximating network, or
whether it generalizes.

\textbf{Density in \(C(K)\) and Stone--Weierstrass.} Let \(K \subset
\mathbb{R}^{d}\) be compact. The Stone--Weierstrass theorem states
that a subalgebra of \(C(K)\) that contains the constants and separates
points is dense in the sup norm. The span of the exponentials \(x \mapsto
e^{w\cdot x + b}\) is such an algebra, since \(e^{a\cdot x}\cdot e^{b\cdot x}
= e^{(a+b)\cdot x}\): networks with exponential units are therefore dense in
\(C(K)\). The exponential is the nonlinearity inside softmax attention.

\textbf{Universality of shallow networks.} For standard activations \(\sigma\)
the span of \(\{\sigma(w\cdot x + b)\}\) is not an algebra, and other arguments are
needed. Cybenko (1989) proved density for sigmoidal activations by
duality: if the span were not dense, Hahn--Banach would give a nonzero
continuous functional vanishing on it, which by the Riesz--Markov
theorem is a nonzero signed measure, and the sigmoidal structure forces
that measure to be zero. Leshno, Lin, Pinkus and Schocken (1993) gave
the sharp statement: for continuous \(\sigma\), the span is dense in \(C(K)\) if and
only if \(\sigma\) is not a polynomial. Density can also be stated in \(L^p(\mu)\);
by \S\ref{sec:math:I4} the two statements are different claims.

\textbf{Rates and regularity classes.} Density gives no rate. Rates
require assuming the target is regular, and the regularity scales are
those of \S\ref{sec:math:I5}.

\begin{itemize}
\tightlist
\item
  \textbf{Sobolev classes and the curse of dimensionality.} For \(f\) in the
  unit ball of \(W^{n,\infty}([0,1]^d)\), Yarotsky (2017)
  shows that deep ReLU networks reach sup-norm accuracy \(\varepsilon\) with
  \(O(\varepsilon^{-d/n} \log(1/\varepsilon))\) weights and \(O(\log(1/\varepsilon))\) depth. The exponent
  \(d/n\) expresses the curse of dimensionality: the required size grows as
  a power of \(1/\varepsilon\) whose exponent is proportional to the dimension and
  inversely proportional to the smoothness.
\item
  \textbf{Optimality.} Continuous nonlinear widths (DeVore, Howard and
  Micchelli) show that \(\varepsilon^{-d/n}\) cannot be beaten when the parameters
  depend continuously on \(f\).~Very deep networks with discontinuous
  parameter selection can beat it (Yarotsky, 2018); this is unrelated to
  the ``super-convergence'' of learning-rate schedules (Smith and Topin).
\item
  \textbf{Besov classes.} Besov spaces capture smoothness that varies
  across the domain, such as functions that are smooth except near local
  singularities. Deep networks adapt to this inhomogeneity, while linear
  methods, which approximate from a fixed subspace, provably cannot
  (Suzuki, 2019).
\item
  \textbf{Barron classes.} For functions with a finite first Fourier
  moment, \(\int |\omega|\, |\hat{f}(\omega)|\, d\omega <
  \infty\), shallow networks with \(N\) units achieve \(L^{2}\) error
  \(O(C_f N^{-1/2})\) (Barron, 1993). The exponent does not depend on \(d\),
  although the constant \(C_f\) can. The proof is
  Fourier-analytic (\S\ref{sec:math:I8}).
\end{itemize}

The hypotheses of these results differ in detail and should be taken
from the original sources; DeVore, Hanin and Petrova, \emph{Neural
Network Approximation} (Acta Numerica, 2021), is a single reference
covering most of them.

\section{The Original Transformer}\label{sec:math-original}

The Transformer of Vaswani et al.,
\emph{\href{https://arxiv.org/abs/1706.03762}{Attention Is All You
Need}} (2017), replaced recurrence and convolution in sequence
transduction with attention alone. This section describes it as the paper
does, with light notation; Section~\ref{sec:math-operators} analyzes it.

\subsection{The task and the overall
architecture}\label{sec:math:II1}

The model was built for machine translation. An encoder maps a source
sequence of symbol representations to continuous representations \(z =
(z_{1}, \ldots, z_{n})\). A decoder then generates the target
sequence one symbol at a time, auto-regressively, consuming its previous
outputs as additional input.

Both stacks have \(N = 6\) identical layers of width \(d_{\text{model}} = 512\).

\begin{itemize}
\tightlist
\item
  \textbf{Encoder layer:} multi-head self-attention, then a
  position-wise feed-forward network.
\item
  \textbf{Decoder layer:} masked multi-head self-attention, then
  multi-head attention over the encoder output, then a feed-forward
  network.
\end{itemize}

Each sub-layer is wrapped in a residual connection followed by layer
normalization, so its output is \(\operatorname{LayerNorm}(x + \operatorname{Sublayer}(x))\). This is the
post-norm arrangement.

\begin{longtable}{@{}>{\raggedright\arraybackslash}p{0.44\linewidth}>{\raggedright\arraybackslash}p{0.44\linewidth}@{}}
\caption{Data flow in the original encoder--decoder Transformer.}\\
\toprule
Stage & Feeds into \\
\midrule
\endfirsthead
\toprule
Stage & Feeds into \\
\midrule
\endhead
\bottomrule
\endlastfoot
Source tokens & Encoder: 6 layers self-attention + FFN \\
Encoder: 6 layers self-attention + FFN & Cross-attention in every decoder layer \\
Previous outputs & Decoder: 6 layers masked self-attention + FFN \\
Cross-attention in every decoder layer & Decoder: 6 layers masked self-attention + FFN \\
Decoder: 6 layers masked self-attention + FFN & Linear + softmax next-token probabilities \\
\end{longtable}

The encoder reads the whole source; the decoder reads its own past
outputs and, through cross-attention, the encoder's output.

\subsection{Scaled dot-product
attention}\label{sec:math:II2}

Queries and keys have dimension \(d_k\) and values dimension \(d_v\). Packing
them into matrices \(Q\), \(K\), \(V\), the attention function is

\[
\operatorname{Attention}(Q, K, V) = \operatorname{softmax}\!\Big(\frac{Q K^\top}{\sqrt{d_k}}\Big) V.
\]

Each output is a weighted sum of values, with weights given by a
compatibility function of the query and the keys. The authors explain
the factor \(1/\sqrt{d_k}\) as a remedy: they suspected that for large \(d_k\) the
dot products grow large and push the softmax into regions with very
small gradients.

\subsection{Multi-head attention}\label{sec:math:II3}

Instead of one attention function at full width, the model projects
queries, keys and values \(h\) times with learned matrices and runs
attention in parallel:

\[
\operatorname{MultiHead}(Q, K, V) = \operatorname{Concat}(\mathrm{head}_1, \dots, \mathrm{head}_h)\, W^O, \qquad \mathrm{head}_i = \operatorname{Attention}(Q W_i^Q, K W_i^K, V W_i^V).
\]

The base model uses \(h = 8\) heads with \(d_k = d_v = 64\). The authors'
motivation is that multiple heads let the model attend jointly to
information from different representation subspaces. In their ablations,
reducing the key size \(d_k\) hurt quality, which they read as a sign that
determining compatibility is not easy.

\subsection{Three uses of
attention}\label{sec:math:II4}

\begin{enumerate}
\def\labelenumi{\arabic{enumi}.}
\tightlist
\item
  \textbf{Encoder self-attention:} queries, keys and values all come
  from the previous encoder layer, and every position attends to every
  position.
\item
  \textbf{Masked decoder self-attention:} each position attends only to
  positions up to and including itself. Illegal connections are masked
  by setting their softmax inputs to \(-\infty\). Combined with shifting the
  output embeddings by one position, this makes the prediction at
  position \(i\) depend only on outputs before \(i\).
\item
  \textbf{Encoder--decoder attention:} queries come from the previous
  decoder layer; keys and values come from the encoder output, so every
  decoder position can attend over the whole input.
\end{enumerate}

\subsection{Feed-forward networks, embeddings, and the output
layer}\label{sec:math:II5}

Each layer applies the same two-layer network to every position
separately, \(\operatorname{FFN}(x) = \max(0, xW_{1} + b_{1})W_{2} +
b_{2}\), with inner dimension 2048. Input and output tokens are
embedded by learned matrices, and a linear map followed by a softmax
converts the decoder output into next-token probabilities. The two
embedding matrices and the pre-softmax linear map share one weight
matrix; in the embedding layers it is multiplied by \(\sqrt{d_{\text{model}}}\).

\subsection{Positional encoding}\label{sec:math:II6}

With no recurrence or convolution, order must be injected explicitly.
The model adds sinusoidal encodings to the input embeddings:

\[
PE_{(pos,\,2i)} = \sin\!\big(pos / 10000^{2i/d_{model}}\big), \qquad PE_{(pos,\,2i+1)} = \cos\!\big(pos / 10000^{2i/d_{model}}\big).
\]

The wavelengths form a geometric progression from \(2\pi\) to \(10000\cdot 2\pi\). The
authors chose sinusoids because, for any fixed offset \(k\), \(PE(pos + k)\) is
a linear function of \(PE(pos)\), which they expected to make attention by
relative position easy to learn.

\subsection{Why self-attention}\label{sec:math:II7}

The paper compares layer types by cost per layer, parallelism, and the
longest path a signal travels between two positions (\(n\) = sequence
length, \(d\) = width, \(k\) = kernel size, \(r\) = neighborhood size).

\begin{longtable}[]{@{}
  >{\raggedright\arraybackslash}p{(\columnwidth - 6\tabcolsep) * \real{0.2500}}
  >{\raggedright\arraybackslash}p{(\columnwidth - 6\tabcolsep) * \real{0.2500}}
  >{\raggedright\arraybackslash}p{(\columnwidth - 6\tabcolsep) * \real{0.2500}}
  >{\raggedright\arraybackslash}p{(\columnwidth - 6\tabcolsep) * \real{0.2500}}@{}}
\toprule\noalign{}
\begin{minipage}[b]{\linewidth}\raggedright
Layer type
\end{minipage} & \begin{minipage}[b]{\linewidth}\raggedright
Complexity per layer
\end{minipage} & \begin{minipage}[b]{\linewidth}\raggedright
Sequential operations
\end{minipage} & \begin{minipage}[b]{\linewidth}\raggedright
Maximum path length
\end{minipage} \\
\midrule\noalign{}
\endhead
\bottomrule\noalign{}
\endlastfoot
Self-attention & \(O(n^{2}\cdot d)\) & \(O(1)\) & \(O(1)\) \\
Recurrent & \(O(n\cdot d^{2})\) & \(O(n)\) & \(O(n)\) \\
Convolutional & \(O(k\cdot n\cdot d^{2})\) & \(O(1)\) & \(O(\log_k(n))\) \\
Restricted self-attention & \(O(r\cdot n\cdot d)\) & \(O(1)\) & \(O(n/r)\) \\
\end{longtable}

Self-attention connects every pair of positions in a constant number of
steps. The cost is quadratic in \(n\), and the authors already proposed
restricting attention to a neighborhood of size \(r\) for very long
sequences.

\subsection{Training and
decoding}\label{sec:math:II8}

\begin{itemize}
\tightlist
\item
  \textbf{Data:} WMT 2014 English--German (about 4.5 million sentence
  pairs, a shared byte-pair vocabulary of about 37,000 tokens) and
  English--French (36 million sentences, a 32,000 word-piece
  vocabulary).
\item
  \textbf{Optimization:} Adam with \(\beta_{1} = 0.9\), \(\beta_{2} = 0.98\) and
  \(\varepsilon = 10^{-9}\), with a learning rate that rises linearly for 4000
  warm-up steps and then decays as the inverse square root of the step
  number.
\item
  \textbf{Regularization:} dropout on sub-layer outputs and on the sums
  of embeddings and positional encodings (rate 0.1 for the base model),
  and label smoothing of 0.1, which the authors note hurts perplexity
  but improves accuracy and BLEU.
\item
  \textbf{Decoding:} beam search with beam size 4 and length penalty
  0.6, with maximum output length set to the input length plus 50 and
  early termination when possible.
\end{itemize}

\subsection{Results}\label{sec:math:II9}

The big model reached 28.4 BLEU on WMT 2014 English--German, more than 2
BLEU above the best previous results, ensembles included, after 3.5 days
on eight P100 GPUs. On English--French it set a single-model state of
the art: 41.0 BLEU in the first arXiv version, reported as 41.8 in the
revised version. A 4-layer model also performed strongly on English
constituency parsing, showing that the architecture generalized beyond
translation.

\section{The Transformer as Operators and Measures}\label{sec:math-operators}

This section rewrites the model of Section~\ref{sec:math-original} in the language of
Section~\ref{sec:math-tools}. The
whole network is a composition of operators on finite-dimensional
Hilbert spaces, written with \(\circ\). Only its last factor has a codomain that
is not a vector space: the set of probability measures on the \(\sigma\)-algebra
of token events, from which the next token is sampled.

\subsection{The Transformer as a composition of
operators}\label{sec:math:IV1}

The 2017 Transformer is a chain of maps, each with a stated domain and
codomain. Every map but the last sends one finite-dimensional Hilbert
space into another; the last sends a Hilbert space into the set of
probability measures on the vocabulary.

\textbf{Notation.} \(\mathcal{V}\) is the vocabulary, with
\(|\mathcal{V}| \approx 37{,}000\) shared byte-pair tokens for
English--German. \(d = d_{\text{model}} = 512\). The source sentence is \(x =
(x_{1}, \ldots, x_n) \in \mathcal{V}^{n}\); the target prefix is
\(y_{<t} = (y_{1}, \ldots, y_{t-1})\). All states are matrices
in \(\mathbb{R}^{n\times d}\), a Hilbert space under the Frobenius inner
product (\S\ref{sec:math:I1}).

\textbf{Building blocks.} The table lists every operator in the order of
evaluation.

\begin{longtable}[]{@{}
  >{\raggedright\arraybackslash}p{(\columnwidth - 6\tabcolsep) * \real{0.2500}}
  >{\raggedright\arraybackslash}p{(\columnwidth - 6\tabcolsep) * \real{0.2500}}
  >{\raggedright\arraybackslash}p{(\columnwidth - 6\tabcolsep) * \real{0.2500}}
  >{\raggedright\arraybackslash}p{(\columnwidth - 6\tabcolsep) * \real{0.2500}}@{}}
\toprule\noalign{}
\begin{minipage}[b]{\linewidth}\raggedright
Operator
\end{minipage} & \begin{minipage}[b]{\linewidth}\raggedright
Map
\end{minipage} & \begin{minipage}[b]{\linewidth}\raggedright
Definition
\end{minipage} & \begin{minipage}[b]{\linewidth}\raggedright
Type
\end{minipage} \\
\midrule\noalign{}
\endhead
\bottomrule\noalign{}
\endlastfoot
One-hot encoding \(\iota\) & \(\mathcal{V}^{n} \to
\mathbb{R}^{n\times |\mathcal{V}|}\) & row \(i\) is the basis vector
\(e_{x_i}\) & injective, measurable \\
Embedding \(\varepsilon\) & \(\mathbb{R}^{n\times |\mathcal{V}|} \to
\mathbb{R}^{n\times d}\) & \(\varepsilon(A) = \sqrt{d}\cdot A W_E\), with \(W_E \in
\mathbb{R}^{|\mathcal{V}|\times d}\) & linear \\
Positional shift \(\tau\) & \(\mathbb{R}^{n\times d} \to
\mathbb{R}^{n\times d}\) & \(\tau(X) = X + P_n\), with \(P_n\) the sinusoidal
table & affine (a translation) \\
Multi-head attention \(M\) & \(\mathbb{R}^{n\times d} \times
\mathbb{R}^{m\times d} \to \mathbb{R}^{n\times d}\) & \(M(X; Y) = \sum_h
\operatorname{softmax}(XW_h^Q \allowbreak (YW_h^K)^{\top} \allowbreak / \sqrt{d_k} + \mathrm{Mask})\allowbreak\, Y W_h^V
W_h^O\) & nonlinear \\
Feed-forward \(F\) & \(\mathbb{R}^{n\times d} \to \mathbb{R}^{n\times d}\)
& \(F(X) = \max(0, XW_{1} + b_{1}) W_{2} + b_{2}\), row by
row & piecewise linear \\
Layer normalization \(\mathrm{LN}\) & \(\mathbb{R}^{n\times d} \to
\mathbb{R}^{n\times d}\) & each row centred, divided by its standard
deviation, then scaled and shifted & nonlinear, smooth \\
Row selection \(\pi_t\) & \(\mathbb{R}^{m\times d} \to \mathbb{R}^{d}\) &
\(\pi_t(Y)\) = row \(t\) of \(Y\) & linear \\
Logit map \(\lambda\) & \(\mathbb{R}^{d} \to
\mathbb{R}^{|\mathcal{V}|}\) & \(\lambda(h) = W_E h\),
the same \(W_E\) as the embedding & linear \\
Softmax \(\sigma\) & \(\mathbb{R}^{|\mathcal{V}|} \to
\Delta(\mathcal{V})\) & \(\sigma(z)_v = e^{z_v} / \sum_u e^{z_u}\) &
nonlinear, into the simplex \\
\end{longtable}

The sum over heads in \(M\) is the concatenation of Section~\ref{sec:math-original} written with
\(W^O\) split into blocks \(W_h^O\). Setting \(Y = X\) gives self-attention;
a mask of \(-\infty\) above the diagonal makes it causal.

\textbf{Sublayers as compositions.} Each sublayer \(G\) is wrapped in a
residual connection and a normalization. As an operator this is a
composition:

\[
R_G \;=\; \mathrm{LN}\circ(\mathrm{Id} + G), \qquad R_G(X) = \mathrm{LN}\big(X + G(X)\big).
\]

\textbf{Encoder.} Encoder layer \(\ell\) applies self-attention
\(S_\ell(X) = M_\ell(X; X)\) and then the feed-forward map
\(F_\ell\). The embedding is \(E = \tau\circ\varepsilon\circ\iota\). The encoder is the six-fold
composition

\[
\mathrm{Enc} \;=\; L^{\mathrm{enc}}_6\circ L^{\mathrm{enc}}_5\circ\cdots\circ L^{\mathrm{enc}}_1\circ E, \qquad L^{\mathrm{enc}}_\ell = R_{F_\ell}\circ R_{S_\ell}.
\]

It maps \(\mathcal{V}^{n}\) to \(\mathbb{R}^{n\times d}\). Its output \(H
= \mathrm{Enc}(x)\) is called the memory.

\textbf{Decoder.} Decoder layer \(\ell\) has three sublayers: causal
self-attention \(\tilde{S}_\ell\), cross-attention \(C_\ell[H] =
M_\ell(Y; H)\) reading the memory, and the feed-forward map. The
memory enters as a fixed parameter of each decoder layer:

\[
L^{\mathrm{dec}}_\ell[H] \;=\; R_{F_\ell}\circ R_{C_\ell[H]}\circ R_{\tilde S_\ell}, \qquad \mathrm{Dec}[H] \;=\; L^{\mathrm{dec}}_6[H]\circ\cdots\circ L^{\mathrm{dec}}_1[H]\circ E .
\]

\textbf{The whole model.} The decoder reads the prefix shifted right by
one position, \((\langle s\rangle, y_{1}, \ldots, y_{t-1})\), where \(\langle s\rangle\) is a start
symbol. The next-token distribution is

\[
p_\theta(\,\cdot \mid x, y_{<t}) \;=\; \big(\sigma\circ\lambda\circ\pi_t\circ \mathrm{Dec}[\mathrm{Enc}(x)]\big)\big(\langle s\rangle, y_1,\dots,y_{t-1}\big).
\]

\(\theta\) collects every weight matrix, bias and normalization parameter.
Written out, one prediction is a composition of \(3 + 6\cdot 2\cdot 2 + 3 + 6\cdot 3\cdot 2 +
3 = 69\) elementary maps (two embeddings of three maps each, every encoder
and decoder sublayer counted as a residual map plus a normalization,
then readout, logits and softmax).

\begin{longtable}{@{}>{\raggedright\arraybackslash}p{0.44\linewidth}>{\raggedright\arraybackslash}p{0.44\linewidth}@{}}
\caption{The Transformer as a chain of operators, from source tokens to a sampled token.}\\
\toprule
Stage & Feeds into \\
\midrule
\endfirsthead
\toprule
Stage & Feeds into \\
\midrule
\endhead
\bottomrule
\endlastfoot
Source \(x \in \mathcal{V}^{n}\) & Embedding \(E\) \\
Embedding \(E\) & Encoder: six layers \\
Encoder: six layers & Memory \(H\) \\
Prefix \(y_{<t} \in \mathcal{V}^{t-1}\) & Embedding \(E\) \\
Embedding \(E\) & Decoder: six layers \\
Memory \(H\) & Decoder: six layers \\
Decoder: six layers & Row \(t\): \(\pi_t\) \\
Row \(t\): \(\pi_t\) & Logits: \(\lambda\) \\
Logits: \(\lambda\) & Softmax: \(\sigma\) \\
Softmax: \(\sigma\) & Measure on \(\mathcal{V}\) \\
Measure on \(\mathcal{V}\) & Sampled token \(y_t\) \\
\end{longtable}

Every factor before \(\sigma\) has a Hilbert space as its codomain. \(\sigma\) changes the
type of object: its values are points of the simplex \(\Delta(\mathcal{V})\),
which \S\ref{sec:math:IV5} identifies with the probability measures on
\((\mathcal{V}, 2^{\mathcal{V}})\). The last arrow is not a
deterministic function of the measure: it requires an independent source
of randomness (\S\ref{sec:math:IV5}).

\subsection{The operators one by
one}\label{sec:math:IV2}

The properties of the composite follow from those of its factors: linear
maps contribute their operator norms, the nonlinear maps their local
behavior, and the chain rule assembles them.

\textbf{Linear and affine factors.} \(\iota\), \(\varepsilon\), \(\pi_t\) and \(\lambda\) are linear, and \(\tau\)
is a translation. Their Lipschitz constants are operator norms (\S\ref{sec:math:I2}):
for instance \(\varepsilon\) is \(\sqrt{d}\cdot\sigma_{\max}(W_E)\)-Lipschitz, where \(\sigma_{\max}\) is the
largest singular value (\S\ref{sec:math:I3}). The input \(x\) itself is discrete;
continuity questions start after \(\iota\), which embeds \(\mathcal{V}^{n}\)
into \(\mathbb{R}^{n\times |\mathcal{V}|}\).

\textbf{Layer normalization is a projection followed by a radial map.}
For a row \(x \in \mathbb{R}^{d}\),

\[
\mathrm{LN}(x) = \gamma \odot \frac{Px}{\sqrt{\|Px\|^2/d + \epsilon}} + \beta, \qquad P = I - \tfrac{1}{d}\mathbf 1\mathbf 1^\top .
\]

\(P\) is the orthogonal projection onto the hyperplane orthogonal to the
all-ones vector (the projection theorem of \S\ref{sec:math:I1}). The radial factor
sends \(Px\) into the ball of radius \(\sqrt{d}\) in that hyperplane. \(\mathrm{LN}\) therefore
maps all of \(\mathbb{R}^{d}\) into a compact set, which is what makes
the finite-precision arguments of \S\ref{sec:math:I9} and \S\ref{sec:math:IV5} possible. \(\mathrm{LN}\) is smooth because \(\epsilon
> 0\). On the hyperplane its Jacobian has norm at most
\(\max|\gamma|/\sqrt{\epsilon}\), so it is globally Lipschitz but with a very
large constant for small \(\epsilon\); for rows with variance near 1 its local
constant is about \(\max|\gamma|\).

\textbf{Feed-forward layers.} \(F = ({\cdot}W_{2} + b_{2})\circ\operatorname{ReLU}\circ
({\cdot}W_{1} + b_{1})\) is piecewise linear. ReLU is 1-Lipschitz, so \(F\)
is \(\|W_{1}\|_{2}\|W_{2}\|_{2}\)-Lipschitz. In the language of
\S\ref{sec:math:I5}, \(F\) belongs to \(W^{1,\infty}_{\mathrm{loc}}\): its weak derivative exists and is
bounded, and its curvature is concentrated on the hyperplanes where a
ReLU unit changes sign.

\textbf{Attention} is real-analytic in its input, hence locally
Lipschitz, but not globally Lipschitz (\S\ref{sec:math:IV3}).

\textbf{Residual maps.} If \(G\) is \(L\)-Lipschitz, \(\mathrm{Id} + G\) is
\((1 + L)\)-Lipschitz. Lipschitz constants multiply under composition, so the
whole chain of \S\ref{sec:math:IV1} is Lipschitz on bounded sets with constant at most
the product of the 69 factors' constants. That bound is valid but far
from sharp; the informative object is the Jacobian at a given input.

\textbf{The chain rule for the whole model.} Write the chain as \(\Phi = \Phi_N
\circ\cdots\circ\Phi_{1}\) with states \(X_k = \Phi_k(X_{k-1})\). Wherever every factor
is differentiable,

\[
D\Phi(X_0) = D\Phi_N(X_{N-1})\, D\Phi_{N-1}(X_{N-2}) \cdots D\Phi_1(X_0), \qquad \nabla_{X_0}\ell = D\Phi(X_0)^\top \nabla \ell .
\]

Backpropagation evaluates the second formula starting from the output:
it applies \(D\Phi_N^{\top}\) to \(\nabla\ell\), then \(D\Phi_{N-1}^{\top}\), and
so on down to the input, multiplying transposed Jacobians onto a vector
without ever forming a full Jacobian matrix. Because ReLU is not
differentiable at zero, some care is needed. Rademacher's theorem
(\S\ref{sec:math:I2}) makes the whole network differentiable almost everywhere,
but it does not by itself guarantee that the product of layer Jacobians
computed by backpropagation equals that derivative: a composition can
sit on a kink of an inner unit on a set of positive measure (a dead unit,
or the convention \(\operatorname{ReLU}'(0) = 0\)). For piecewise-analytic networks
such as this one, backpropagation returns the true gradient almost
everywhere and an element of a conservative field everywhere (Kakade and
Lee, 2018; Bolte and Pauwels, 2021).

\textbf{Post-norm changes the Jacobian of each block.} For \(R_G = \mathrm{LN}\circ
(\mathrm{Id} + G)\),

\[
DR_G(X) = D\mathrm{LN}\big(X + G(X)\big)\,\big(I + DG(X)\big).
\]

The identity term is multiplied by the Jacobian of the normalization. A
gradient reaching the encoder input passes through the normalization
Jacobian of every sublayer on its path, and there is no identity path
through the stack. Xiong et al.~(2020) relate this to the 2017 model's
dependence on learning-rate warm-up; the pre-norm form \(X + G(\mathrm{LN}(X))\),
whose Jacobian is \(I + DG\cdot D\mathrm{LN}\), restores a direct identity path.

\textbf{Summary of regularity.} The composite map from the one-hot input
to the logits is piecewise real-analytic: analytic on each region where
the ReLU activation pattern is fixed, locally Lipschitz everywhere,
differentiable almost everywhere, and in \(W^{1,\infty}_{\mathrm{loc}}\). The softmax
that follows is analytic, so every next-token probability is a locally
Lipschitz function of the embedded input and of \(\theta\).

\subsection{Attention as a data-dependent Markov
kernel}\label{sec:math:IV3}

Inside the network, probability measures already appear before the last
layer: every attention head builds one measure over positions for each
query, and outputs an expectation under it.

\textbf{One head.} For queries from \(X \in \mathbb{R}^{n\times d}\) and
keys and values from \(Y \in \mathbb{R}^{m\times d}\), the head computes
scores, a row-stochastic matrix, and values:

\[
A = \frac{XW^Q\,(YW^K)^\top}{\sqrt{d_k}} + \mathrm{Mask}, \qquad P = \operatorname{softmax}_{\mathrm{row}}(A), \qquad \mathrm{head}(X;Y) = P\,YW^V .
\]

Row \(i\) of \(P\) is a probability measure on the finite set of key positions
\(\{1, \ldots, m\}\) with its power-set \(\sigma\)-algebra: a Markov kernel from
query positions to key positions (\S\ref{sec:math:I6}). Row \(i\) of the output is an
expectation:

\[
\mathrm{head}(X;Y)_i = \sum_{j=1}^{m} P_{ij}\,(YW^V)_j = \mathbb E_{j\sim P_{i\cdot}}\big[(YW^V)_j\big].
\]

The kernel is recomputed from the input at every forward pass. This
dependence is the source of attention's nonlinearity: for fixed \(P\) the
map \(Y \mapsto P Y W^V\) would be linear.

\textbf{Three kernels in the 2017 model.}

\begin{longtable}[]{@{}
  >{\raggedright\arraybackslash}p{(\columnwidth - 6\tabcolsep) * \real{0.2500}}
  >{\raggedright\arraybackslash}p{(\columnwidth - 6\tabcolsep) * \real{0.2500}}
  >{\raggedright\arraybackslash}p{(\columnwidth - 6\tabcolsep) * \real{0.2500}}
  >{\raggedright\arraybackslash}p{(\columnwidth - 6\tabcolsep) * \real{0.2500}}@{}}
\toprule\noalign{}
\begin{minipage}[b]{\linewidth}\raggedright
Sublayer
\end{minipage} & \begin{minipage}[b]{\linewidth}\raggedright
Kernel
\end{minipage} & \begin{minipage}[b]{\linewidth}\raggedright
Shape
\end{minipage} & \begin{minipage}[b]{\linewidth}\raggedright
Support of row \(i\)
\end{minipage} \\
\midrule\noalign{}
\endhead
\bottomrule\noalign{}
\endlastfoot
Encoder self-attention & source positions to source positions & \(n \times n\),
all entries positive & all of \(\{1, \ldots, n\}\) \\
Decoder self-attention & target positions to target positions & \(m \times m\),
lower triangular & \(\{1, \ldots, i\}\) \\
Cross-attention & target positions to source positions & \(m \times n\), all
entries positive & all of \(\{1, \ldots, n\}\) \\
\end{longtable}

A causal mask sets scores above the diagonal to \(-\infty\). By \S\ref{sec:math:I7} this is
conditioning: row \(i\) is the unmasked measure conditioned on the event
\(\{j \le i\}\).

\textbf{Norms.} Each output coordinate is a convex combination of value
coordinates, so \(P\) is non-expansive in the maximum norm: \(\|P\|_{\infty\to\infty} = 1\).
In the Euclidean norm it is not: \(\|P\|_{2}\) can reach \(\sqrt{m}\) (\S\ref{sec:math:I3}).

\textbf{Local but not global Lipschitz continuity.} Each row of the
softmax has Jacobian \(\operatorname{diag}(p) - pp^{\top}\) (\S\ref{sec:math:I7}), and the scores are
quadratic in the input. The head is therefore real-analytic and locally
Lipschitz, but its sensitivity grows with the size of its input, and
dot-product self-attention is not globally Lipschitz (Kim, Papamakarios
and Mnih, 2021). In the post-norm
2017 model every attention sublayer after the first reads a
layer-normalized state, which lies in a compact set (\S\ref{sec:math:IV2}). On that set
attention is Lipschitz. The factor \(1/\sqrt{d_k}\) keeps the score variance near
1 when query and key components have unit variance, which keeps the
softmax away from saturation.

\textbf{Permutation equivariance of the encoder.} Let \(\Pi\) be an \(n \times n\)
permutation matrix. With no mask, \(\operatorname{softmax}_{\mathrm{row}}(\Pi A\Pi^{\top}) = \Pi\,
\operatorname{softmax}_{\mathrm{row}}(A)\, \Pi^{\top}\), and therefore

\[
S(\Pi X) = \Pi\, S(X), \qquad F(\Pi X) = \Pi F(X), \qquad \mathrm{LN}(\Pi X) = \Pi\,\mathrm{LN}(X).
\]

Every encoder layer is permutation-equivariant. Only the positional
translation \(\tau\) breaks the symmetry, since \(P_n\) is not
permutation-invariant: all information about word order in the encoder
comes from \(\tau\). The decoder's causal mask is not permutation-invariant, so
the decoder is not equivariant, and the mask itself carries order
information.

\textbf{Cross-attention depends on the memory only through its empirical
measure.} Permuting the rows of the memory \(H\) permutes keys and values
together, so \(C[\Pi H] = C[H]\). The output depends on \(H\) only through
the empirical measure \(\mu_H = (1/n) \sum_j \delta_{h_j}\) of its rows:

\[
C[H](y) = \frac{\int e^{\langle yW^Q,\, hW^K\rangle/\sqrt{d_k}}\; hW^V\, d\mu_H(h)}{\int e^{\langle yW^Q,\, hW^K\rangle/\sqrt{d_k}}\, d\mu_H(h)} .
\]

This is the Nadaraya--Watson form of the integral operator of \S\ref{sec:math:I4} with
the exponential kernel. The positive definiteness of \S\ref{sec:math:I4}, and with it
Mercer's theorem, holds in the projected variables \(q = yW^Q\) and \(k = hW^K\).
As a function of \((y, h)\) itself the kernel is \(e^{yBh^{\top}/\sqrt{d_k}}\) with
\(B = W^Q (W^K)^\top\), which is neither symmetric nor positive definite. Order in the source survives only because \(\tau\) has
written positions into the rows of \(H\).

\textbf{Stacking kernels.} Composing attention layers multiplies
row-stochastic matrices, as in a Markov chain. The Dobrushin coefficient
\(\delta(P) = \tfrac12 \max_{i,k} \sum_j |P_{ij} - P_{kj}|\) is
submultiplicative, \(\delta(PQ) \le \delta(P)\, \delta(Q)\), and it measures how far apart the
rows of \(P\) are. Every entry of an unmasked softmax kernel is positive, so
\(\delta(P) < 1\): each unmasked layer strictly contracts the
differences between rows. Pure attention stacks therefore drive all rows
toward a common value and lose rank with depth (Dong, Cordonnier and
Loukas, 2021). The residual maps \(\mathrm{Id} + G\) and the feed-forward layers of
\S\ref{sec:math:IV1} are what prevent this collapse.

\subsection{Spectral structure: heads, weight tying and positional
rotations}\label{sec:math:IV4}

The nonlinear operators of \S\ref{sec:math:IV1} are built from linear maps, and those
linear maps have the spectral structure of \S\ref{sec:math:I3}.

\textbf{Each head compares tokens through a low-rank bilinear form.}
Before the softmax, head \(h\) scores a query row \(x\) against a key row \(y\) by

\[
b_h(x, y) = \frac{x^\top B_h\, y}{\sqrt{d_k}}, \qquad B_h = W_h^Q (W_h^K)^\top \in \mathbb R^{512\times 512}, \qquad \operatorname{rank} B_h \le d_k = 64 .
\]

With the SVD \(B_h = \sum_{r\le 64} s_r u_r v_r^{\top}\), the head
compares the components of the query along \(u_r\) with those of the key
along \(v_r\), weighted by \(s_r\). It depends only on the projections of the
query and the key onto 64-dimensional subspaces of the 512-dimensional
representation.

\textbf{Asymmetry.} Split \(B_h\) into symmetric and skew-symmetric parts,
\(B_h = \tfrac12(B_h + B_h^{\top}) + \tfrac12(B_h - B_h^{\top})\). The skew
part is why \(b_h(x, y) \neq b_h(y, x)\) in general: token \(i\) can assign large
weight to token \(j\) while \(j\) assigns small weight to \(i\). The skew part is not
the only source: the softmax normalizes each row separately, so the
attention weights are asymmetric in general even when \(B_h\) is symmetric.

\textbf{The output of each head lies in a low-rank subspace.} The value
and output maps combine into \(C_h = W_h^V W_h^O \in
\mathbb{R}^{512\times 512}\), also of rank at most 64. Multi-head
attention is

\[
M(X;Y) = \sum_{h=1}^{8} P_h(X;Y)\; Y\, C_h ,
\]

a sum of eight terms, each confined to the at most 64-dimensional row
space of its \(C_h\). This is the precise form of the paper's motivation
that heads attend to information from different representation
subspaces. By Eckart--Young (\S\ref{sec:math:I3}), the singular values of \(B_h\) and \(C_h\)
measure how far each head can be compressed.

\textbf{Scores are low rank; the kernel need not be.} The score matrix
of a head has rank at most 64, but the exponential inside the softmax
acts entrywise and can raise the rank. The kernel \(P_h\) may be full rank
even though its scores are not.

\textbf{Weight tying and Riesz representation.} The logit map uses the
embedding matrix: \(\lambda(h)_v = \langle e_v, h\rangle\), where \(e_v\) is row \(v\) of \(W_E\). By
the Riesz representation theorem (\S\ref{sec:math:I1}), every linear functional on
\(\mathbb{R}^{d}\) is an inner product with a unique vector. Tying
weights means that the vector representing token \(v\) at the input is also
the representer of the functional that scores \(v\) at the output.

\textbf{Positional encoding is a rotation.} Group the 512 coordinates of
\(P_n\) into 256 pairs, pair \(j = 0, \ldots, 255\) oscillating at frequency \(\omega_j =
10000^{-2j/512}\). Write pair \(j\) as a complex number, \(z_j(\mathrm{pos}) =
e^{i\omega_j \cdot \mathrm{pos}}\), with \(i\) the imaginary unit. A shift by \(k\) positions is
then multiplication by \(e^{i\omega_j k}\):

\[
PE(pos + k) = R(k)\, PE(pos), \qquad R(k) = \operatorname{diag}\big(R_{-\omega_0 k}, \dots, R_{-\omega_{255} k}\big),
\]

with \(R_\theta\) the \(2 \times 2\) rotation by \(\theta\). The sign is negative because each
pair is stored in the order \((\sin, \cos)\). \(R(k)\) is orthogonal, so all its
singular values are 1: this is the linear shift property the authors
cite, and the shift theorem of \S\ref{sec:math:I8}. A single pair determines position
only modulo its period \(2\pi/\omega_j\); the geometric progression of periods,
from \(2\pi\) to \(2\pi \cdot 10000^{510/512} \approx 60{,}611\), resolves positions over a
range set by the longest period (\S\ref{sec:math:I8}).

\textbf{Where position enters the score.} Because \(\tau\) adds \(P_n\) to the
content, a score expands into four terms:

\[
(x_i + p_i)^\top B\,(x_j + p_j) = x_i^\top B x_j + x_i^\top B p_j + p_i^\top B x_j + p_i^\top B p_j .
\]

The last term depends on positions only, but in general on both \(i\) and \(j\),
not only on \(i - j\): \(p_i^{\top} B R(j - i) p_i\) still depends on \(p_i\)
unless \(B\) commutes with the rotations. Relative position is available to
the model but not built into the score. Later rotary encodings (Su et al., 2021) multiply
queries and keys by rotations instead of adding, which makes every score
depend on position only through \(i - j\).

\subsection{\texorpdfstring{The last layer: a probability measure on the \(\sigma\)-algebra
of token events}{The last layer: a probability measure on the sigma-algebra of token events}}\label{sec:math:IV5}

The last three operators of \S\ref{sec:math:IV1}, \(\pi_t\), \(\lambda\) and \(\sigma\), turn a vector of the
Hilbert space \(\mathbb{R}^{d}\) into a probability measure on a finite
measurable space. Sampling then turns that measure into a token. This
section makes each step precise.

\textbf{Step 1: logits.} The decoder's last state at position \(t\) is \(h_t
= \pi_t(\mathrm{Dec}[H](\langle\mathrm{s}\rangle, y_{1}, \ldots, y_{t-1})) \in
\mathbb{R}^{512}\). The logit map gives one real number per token:

\[
z = \lambda(h_t) = W_E\, h_t \in \mathbb R^{|\mathcal V|}, \qquad z_v = \langle e_v, h_t\rangle .
\]

Each logit is a linear functional of \(h_t\) with Riesz representer \(e_v\)
(\S\ref{sec:math:IV4}). Adding the same constant to every logit changes nothing below,
so logits are really coordinates in the quotient space
\(\mathbb{R}^{|\mathcal{V}|} /
\mathbb{R}\mathbf{1}\) (\S\ref{sec:math:I7}).

\textbf{Step 2: the measurable space of token events.} The next token
takes values in the sample space \(\mathcal{V}\). Because
\(\mathcal{V}\) is finite, the natural \(\sigma\)-algebra is its power set:

\[
(\mathcal V,\ 2^{\mathcal V}), \qquad |2^{\mathcal V}| = 2^{|\mathcal V|} \approx 2^{37\,000}.
\]

Every subset of the vocabulary is an event, closed under complement and
union as \S\ref{sec:math:I6} requires. The atoms are the singletons \(\{v\}\), and every
event is a finite union of atoms. Events have direct linguistic
readings:

\begin{itemize}
\tightlist
\item
  \(\{v : v \text{ is a punctuation token}\}\), ``the next token is punctuation'';
\item
  \(\{v : v \text{ begins a new word}\}\), read off the byte-pair markers;
\item
  \(\{y_t\}\), ``the model produces the reference token'';
\item
  the top-\(k\) set of the current logits, an event chosen by the logits
  themselves.
\end{itemize}

No measurability issues arise on this side: on a finite set every
function is measurable. The continuous states of the network form the
domain of the kernel, \(\mathbb{R}^{d}\) with its Borel \(\sigma\)-algebra.

\textbf{Step 3: the probability measure.} The softmax assigns a mass to
each atom, and additivity extends it to every event:

\[
P_t(S) \;=\; \sum_{v\in S} \sigma(z)_v \;=\; \sum_{v\in S}\frac{e^{z_v}}{\sum_{u\in\mathcal V} e^{z_u}}, \qquad S\in 2^{\mathcal V}.
\]

The axioms hold directly: \(P_t(S) \ge 0\), \(P_t(\mathcal{V}) = 1\), and
\(P_t(S \cup S') = P_t(S) + P_t(S')\) for disjoint \(S\) and \(S'\). Since the
\(\sigma\)-algebra is finite, countable additivity reduces to finite additivity.
Every atom has positive mass, so the only null event is \(\emptyset\): \(P_t\) has the
same null sets as the counting measure.

\textbf{The simplex is the space of measures.} A probability measure on
\((\mathcal{V}, 2^{\mathcal{V}})\) is determined by its values on
atoms. The map \(p \mapsto \sum_v p_v \delta_v\) identifies the simplex
\(\Delta(\mathcal{V}) = \{p \in
\mathbb{R}^{|\mathcal{V}|} : p \ge 0, \sum_v p_v
= 1\}\) with the set of all probability measures on \((\mathcal{V},
2^{\mathcal{V}})\). The softmax is a bijection from
\(\mathbb{R}^{|\mathcal{V}|}/\mathbb{R}\mathbf{1}\)
onto the open simplex, the measures that charge every atom; its inverse
is \(p \mapsto \log p\) up to a constant. The logits are thus the natural
parameters of the categorical exponential family, and \(\sigma\) is the operator
of \S\ref{sec:math:IV1} that changes the type of object: its codomain is a convex set
of measures, not a vector space.

\textbf{Step 4: the whole model is a Markov kernel.} Write \(\Phi_\theta\) for the
deterministic part, the map from the context \((x, y_{<t})\) to
\(h_t\), and \(\kappa_\theta(h, S) = \sum_{v\in S} \sigma(W_E h)_v\) for the last layer. \(\kappa_\theta\) is
a Markov kernel from \((\mathbb{R}^{d}, \mathrm{Borel})\) to \((\mathcal{V},
2^{\mathcal{V}})\): for each \(h\) it is a probability measure, and for
each event \(S\) the map \(h \mapsto \kappa_\theta(h, S)\) is continuous, hence Borel
measurable (\S\ref{sec:math:I6}). The model's conditional law is the kernel composed
with the operator chain:

\[
P_\theta(S \mid x, y_{<t}) \;=\; \kappa_\theta\big(\Phi_\theta(x, y_{<t}),\, S\big), \qquad \Phi_\theta = \pi_t\circ\mathrm{Dec}[\mathrm{Enc}(x)] .
\]

Every factor of \(\Phi_\theta\) is a deterministic map between Hilbert spaces;
randomness enters only through \(\kappa_\theta\).

\textbf{Operations on the measure.} Decoding procedures act on \(P_t\)
through the structure of \S\ref{sec:math:I7}:

\begin{longtable}[]{@{}
  >{\raggedright\arraybackslash}p{(\columnwidth - 4\tabcolsep) * \real{0.3333}}
  >{\raggedright\arraybackslash}p{(\columnwidth - 4\tabcolsep) * \real{0.3333}}
  >{\raggedright\arraybackslash}p{(\columnwidth - 4\tabcolsep) * \real{0.3333}}@{}}
\toprule\noalign{}
\begin{minipage}[b]{\linewidth}\raggedright
Operation
\end{minipage} & \begin{minipage}[b]{\linewidth}\raggedright
On the logits
\end{minipage} & \begin{minipage}[b]{\linewidth}\raggedright
On the measure
\end{minipage} \\
\midrule\noalign{}
\endhead
\bottomrule\noalign{}
\endlastfoot
Temperature \(\tau\) & \(z \mapsto z/\tau\) & moves along the Gibbs family; the maximizer of
\(\langle p, z\rangle + \tau H(p)\) \\
\(\tau \to 0\) & --- & \(P_t \to \delta\) at the most likely token, when it is unique \\
\(\tau \to \infty\) & --- & \(P_t \to\) the uniform measure on \(\mathcal{V}\) \\
Top-\(k\) or nucleus restriction to \(S\) & \(z_v \mapsto -\infty\) for \(v \notin S\) & conditioning:
\(P_t(\cdot \mid S) = P_t(\cdot \cap S) / P_t(S)\) \\
\end{longtable}

\textbf{Step 5: sampling a token.} A measure is not a token. Sampling
means constructing a random variable whose law is \(P_t\). Take a
probability space carrying one uniform variable, (\([0, 1]\), Borel
sets, Lebesgue measure), order the vocabulary as \(v_{1}, \ldots,
v_K\), and let \(F_k = \sum_{j\le k} P_t(\{v_j\})\) with \(F_{0} = 0\). Define

\[
Y_t(u) = v_k \quad\text{when}\quad F_{k-1} < u \le F_k .
\]

\(Y_t\) is measurable, because the preimage of each atom \(\{v_k\}\) is the
interval \((F_{k-1}, F_k]\). Its law, the pushforward of Lebesgue
measure, is exactly \(P_t\):

\[
\mathrm{Leb}\big(Y_t^{-1}(\{v_k\})\big) = F_k - F_{k-1} = P_t(\{v_k\}).
\]

The Gumbel-max construction gives the same law without normalizing: with
\(G_v\) independent standard Gumbel variables, \(\arg\max_v (z_v + G_v)\) has
law \(\sigma(z)\) (\S\ref{sec:math:I7}). In both constructions the randomness is external and
independent of \(\theta\): the model supplies the measure, and a separate source
of noise supplies the draw.

\textbf{What training asks of the measure.} During training the
reference token \(y_t\) defines a target measure. With label smoothing \(\varepsilon =
0.1\), it is

\[
q_t = (1-\varepsilon)\,\delta_{y_t} + \varepsilon\, u_{\mathcal V}, \qquad \mathcal L_t = H(q_t, P_t) = H(q_t) + \mathrm{KL}(q_t \,\|\, P_t),
\]

with \(u_{\mathcal{V}}\) the uniform measure. (Implementations differ in detail; some
spread \(\varepsilon\) over the \(|\mathcal{V}|-1\) non-reference tokens only. The argument
below is unchanged.) By Gibbs' inequality
(\S\ref{sec:math:I10}) the loss is minimized at \(P_t = q_t\), a measure that charges
every atom; smoothing thus keeps the target inside the open simplex,
which is the range of the softmax. How training selects \(\theta\) is outside the
scope of this chapter.

\textbf{In floating point.} The exact \(P_t\) charges every atom, but
computed probabilities can underflow to zero, so the computed measure
can have null events the ideal one lacks. Implementations compute \(\log
\sigma(z) = z - \max z - \log \sum_u e^{z_u - \max z}\), using the shift
invariance of Step 1 to avoid overflow.

\subsection{Sequences: from one token to a measure on
translations}\label{sec:math:IV6}

Section IV.5 produced one measure per position. Chained along the target
sentence, these measures define a single probability measure on the
space of whole translations, conditional on the source.

\textbf{The sample space of translations.} Fix a source sentence \(x\).
Outputs are sequences \(\omega = (y_{1}, y_{2}, \ldots)\) in
\(\mathcal{V}^{\mathbb{N}}\); a finite translation is a sequence
that reaches the end symbol \(\langle/\mathrm{s}\rangle\) and stays there. The \(\sigma\)-algebra is
generated by the cylinder sets of \S\ref{sec:math:I6},

\[
C(y_1,\dots,y_t) = \{\omega\in\mathcal V^{\mathbb N} : \omega_1 = y_1,\dots,\omega_t = y_t\}, \qquad \mathcal F_t = \sigma(\omega_1,\dots,\omega_t),
\]

and the filtration \(\mathcal{F}_{0} \subset \mathcal{F}_{1} \subset \mathcal{F}_{2} \subset \cdots\) records what is
known after \(t\) tokens. The event ``the translation begins with Das Haus''
is a cylinder of length 2 and belongs to \(\mathcal{F}_{2}\).

\textbf{Predictability: the causal mask as a measurability statement.}
The decoder input at position \(s\) is \(y_{s-1}\), shifted right by one. By
induction over the layers of \S\ref{sec:math:IV1}, row \(t\) of every decoder layer depends
only on the inputs at positions \(1, \ldots, t\) and on the memory \(H\):

\begin{itemize}
\tightlist
\item
  causal self-attention: row \(t\) of its kernel is supported on \(\{1,
  \ldots, t\}\) (\S\ref{sec:math:IV3});
\item
  cross-attention: it reads \(H = \mathrm{Enc}(x)\), which depends on \(x\) alone;
\item
  feed-forward layers, normalization and residual maps act row by row.
\end{itemize}

So \(h_t\) is a function of \((x, y_{1}, \ldots, y_{t-1})\). By the
Doob--Dynkin lemma (\S\ref{sec:math:I6}), the next-token kernel is \(\mathcal{F}_{t-1}\)-measurable
for fixed \(x\):

\[
y \mapsto p_\theta(\,\cdot\mid x, y_{<t}) \ \text{is}\ \mathcal F_{t-1}\text{-measurable for each fixed } x.
\]

The sequence of kernels is predictable: each distribution is fixed
before its token is drawn. The encoder's output, by contrast, depends on the whole source at once.
The decoder's causal mask and the encoder's unmasked self-attention are
thus two measurability conditions.

\textbf{One measure on all translations.} By the Ionescu-Tulcea theorem
(\S\ref{sec:math:I6}), the predictable kernels determine a unique probability measure
\(P_\theta(\cdot \mid x)\) on \((\mathcal{V}^{\mathbb{N}}, \mathcal{F}_\infty)\). On cylinders
it is given by the chain rule:

\[
P_\theta\big(C(y_1,\dots,y_t)\mid x\big) = \prod_{s=1}^{t} p_\theta(y_s \mid x, y_{<s}).
\]

The cylinders form a \(\pi\)-system generating \(\mathcal{F}_\infty\), so by the \(\pi\)-\(\lambda\) theorem
these values determine \(P_\theta(\cdot \mid x)\) uniquely. Letting \(x\) vary, \(x \mapsto P_\theta(\cdot
\mid x)\) is a Markov kernel from source sentences to target sequences: this
kernel is the translation model.

\textbf{Generation realizes the measure path by path.} Draw independent
uniforms \(U_{1}, U_{2}, \ldots\) and set \(Y_t = G(p_\theta(\cdot \mid x,
Y_{<t}), U_t)\), with \(G\) the inverse-distribution construction
of \S\ref{sec:math:IV5}. Each \(Y_t\) is \(\sigma(U_1, \ldots, U_t)\)-measurable, and the law of \((Y_{1},
Y_{2}, \ldots)\) is exactly \(P_\theta(\cdot \mid x)\). Equivalently, generation is a
Markov chain on the countable set of prefixes \(\mathcal{V}^*\), whose
transitions append one token until \(\langle/\mathrm{s}\rangle\) is emitted. States ending in
\(\langle/\mathrm{s}\rangle\) are absorbing: once the end symbol appears the chain stays
put, which amounts to a cemetery state that emits nothing further.

\textbf{Termination.} The first time the end symbol appears, \(T = \inf\{t
: Y_t = \langle/\mathrm{s}\rangle\}\), is a stopping time: \(\{T = t\} \in \mathcal{F}_t\). Generation ends
almost surely exactly when \(P_\theta(T < \infty \mid x) = 1\), the tightness
condition; in Markov-chain terms, the absorbing state is reached with
probability one. The 2017 decoder also capped output length at the input
length plus 50, which makes termination certain in practice.

\textbf{The target of the kernel.} Let \(P^*\) be the true joint law of
source and target sentences. The ideal next-token law is a conditional
expectation, and by \S\ref{sec:math:I6} an orthogonal projection:

\[
P^*(Y_t = v \mid x, Y_{<t}) = \mathbb E^*\big[\mathbf 1_{\{Y_t = v\}} \,\big|\, \sigma(x)\vee\mathcal F_{t-1}\big],
\]

the projection of the indicator of the event \(\{Y_t = v\}\) onto the
subspace of \(L^{2}(P^*)\) of functions of the source and the prefix. By
Doob--Dynkin it is a function of \((x, y_{<t})\), and the decoder is
a parametrized family of exactly such functions. The divergence between
the true and the model measures on whole sequences decomposes along the
filtration as a sum of per-token divergences (\S\ref{sec:math:I10}), which is why a
per-token loss trains a sequence model.

\subsection{Decoding and the idealized
model}\label{sec:math:IV7}

The measure \(P_\theta(\cdot \mid x)\) is the model; a decoding rule is a separate
procedure applied to it. The 2017 paper decoded by search rather than by
sampling.

\textbf{Sampling versus search.} Three rules act on the same measure:

\begin{longtable}[]{@{}
  >{\raggedright\arraybackslash}p{(\columnwidth - 4\tabcolsep) * \real{0.3333}}
  >{\raggedright\arraybackslash}p{(\columnwidth - 4\tabcolsep) * \real{0.3333}}
  >{\raggedright\arraybackslash}p{(\columnwidth - 4\tabcolsep) * \real{0.3333}}@{}}
\toprule\noalign{}
\begin{minipage}[b]{\linewidth}\raggedright
Rule
\end{minipage} & \begin{minipage}[b]{\linewidth}\raggedright
Output
\end{minipage} & \begin{minipage}[b]{\linewidth}\raggedright
Relation to \(P_\theta(\cdot \mid x)\)
\end{minipage} \\
\midrule\noalign{}
\endhead
\bottomrule\noalign{}
\endlastfoot
Ancestral sampling (\S\ref{sec:math:IV6}) & a random translation & a draw whose law is
exactly \(P_\theta(\cdot \mid x)\) \\
Greedy decoding & \(\arg\max_v p_\theta(v \mid x, y_{<t})\) at each step &
the \(\tau \to 0\) limit of per-step sampling; not in general the most likely
sequence \\
Beam search, beam 4, length penalty \(\alpha = 0.6\) & an approximate maximizer
of \(\log P_\theta(y \mid x) / \mathrm{lp}(y)\) & a deterministic functional of the measure,
not a draw from it \\
\end{longtable}

Beam search keeps the four best prefixes at each step and compares
finished candidates after dividing their log-probability by a length
penalty \(\mathrm{lp}(y)\) that grows with length, so that shorter outputs are not
automatically preferred. It is a heuristic for the mode of a reweighted
sequence measure. The BLEU scores of Section~\ref{sec:math-original} therefore measure a
mode-seeking functional of \(P_\theta\), not typical samples from it.

\textbf{The model that was computed.} The analysis of Parts I and IV
concerns the exact operators. Three features of the 2017 setup separate
the computed model from that idealization:

\begin{itemize}
\tightlist
\item
  \textbf{Finite precision.} Every operator was evaluated in floating
  point on P100 GPUs, so the states lie on a finite grid and small
  probabilities can underflow (\S\ref{sec:math:I9}, \S\ref{sec:math:IV5}).
\item
  \textbf{Dropout.} During training, dropout (rate 0.1 for the base model, 0.3 for the big English--German
  model) multiplies
  sublayer outputs by random masks. The training-time network is
  therefore itself a random operator, a Markov kernel from inputs to
  states, and only the test-time network, with dropout off, is the
  deterministic composition of \S\ref{sec:math:IV1}.
\item
  \textbf{Finite length.} The decoder never runs beyond input length
  plus 50 tokens, so the measure actually realized is the image of
  \(P_\theta(\cdot \mid x)\) under truncation at that length.
\end{itemize}

\textbf{The chain as a whole.} Each stage of the model lives on a
different mathematical object, and each passage between stages is one of
the operators or constructions of Section~\ref{sec:math-tools}:

\begin{longtable}[]{@{}
  >{\raggedright\arraybackslash}p{(\columnwidth - 6\tabcolsep) * \real{0.2500}}
  >{\raggedright\arraybackslash}p{(\columnwidth - 6\tabcolsep) * \real{0.2500}}
  >{\raggedright\arraybackslash}p{(\columnwidth - 6\tabcolsep) * \real{0.2500}}
  >{\raggedright\arraybackslash}p{(\columnwidth - 6\tabcolsep) * \real{0.2500}}@{}}
\toprule\noalign{}
\begin{minipage}[b]{\linewidth}\raggedright
Stage
\end{minipage} & \begin{minipage}[b]{\linewidth}\raggedright
Object
\end{minipage} & \begin{minipage}[b]{\linewidth}\raggedright
Space
\end{minipage} & \begin{minipage}[b]{\linewidth}\raggedright
Structure used
\end{minipage} \\
\midrule\noalign{}
\endhead
\bottomrule\noalign{}
\endlastfoot
Input & source \(x\), prefix \(y_{<t}\) & \(\mathcal{V}^{n}\),
\(\mathcal{V}^{t- 1}\) & finite sets, power-set \(\sigma\)-algebras \\
Embedding & \(E = \tau \circ \varepsilon \circ \iota\) & \(\mathbb{R}^{n\times d}\) & affine map into
a Hilbert space \\
Encoder and decoder & 12 layers composed with \(\circ\) &
\(\mathbb{R}^{n\times d}\), \(\mathbb{R}^{m\times d}\) & nonlinear
operators; attention kernels inside \\
Readout & \(\lambda \circ \pi_t\) &
\(\mathbb{R}^{|\mathcal{V}|} /
\mathbb{R}\mathbf{1}\) & linear functionals, Riesz representers \\
Softmax & \(\sigma\) & \(\Delta(\mathcal{V})\) & probability measures on
\((\mathcal{V}, 2^{\mathcal{V}})\) \\
Sampling & \(Y_t\) & \(\mathcal{V}\) & a random variable with law \(P_t\) \\
Sequence & \(P_\theta(\cdot \mid x)\) & \((\mathcal{V}^{\mathbb{N}}, \mathcal{F}_\infty)\) &
filtration, predictable kernels, Ionescu-Tulcea \\
\end{longtable}

Read from top to bottom, this is the thesis of the chapter: the
Transformer is a composition of operators on Hilbert spaces, followed by
a single map into the space of probability measures on the \(\sigma\)-algebra of
token events; sampling, applied step after step along a filtration,
turns that measure into text.

\section{Beyond the Structure of the Family}

A transformer, as described in this chapter, is a map \(\theta \mapsto P_\theta\) from
parameters to measures on sequence space. The geometry of Hilbert spaces
carries its representations; spectral theory governs its parameters and
local behavior; measure-theoretic probability, organized by the
\(\sigma\)-algebra of events and its filtration, turns its output into
distributions over tokens and sequences. Everything here concerns the
structure of that family. The natural next question, which this chapter
leaves open, is which member to choose:
pretraining as maximum likelihood, fine-tuning, reinforcement learning
with a KL penalty, whose optimum is a Gibbs tilt of a reference measure,
and on-policy distillation between sequence measures.

\section{Sources}

Primary sources:

\begin{itemize}
\tightlist
\item
  Vaswani et al., \href{https://arxiv.org/abs/1706.03762}{Attention Is
  All You Need}, arXiv:1706.03762 (first version June 2017; results
  quoted from v1 unless noted).
\end{itemize}

Further references:

\begin{itemize}
\tightlist
\item
  Approximation theory: Cybenko (1989); Leshno, Lin, Pinkus and Schocken
  (1993); Barron (1993); Yarotsky (2017, 2018); DeVore, Howard and
  Micchelli; Suzuki (2019); DeVore, Hanin and Petrova, \emph{Neural
  Network Approximation}, Acta Numerica (2021).
\item
  Layer normalization placement: Xiong et al., \emph{On Layer
  Normalization in the Transformer Architecture}, ICML (2020).
\item
  Rank collapse in attention stacks: Dong, Cordonnier and Loukas,
  \emph{Attention is not all you need: pure attention loses rank doubly
  exponentially with depth}, ICML (2021).
\item
  Kernel methods: Steinwart and Christmann, \emph{Support Vector
  Machines}, for Mercer's theorem and eigenvalue decay.
\item
  Lipschitz continuity of attention: Kim, Papamakarios and Mnih,
  \emph{The Lipschitz Constant of Self-Attention}, ICML (2021).
\item
  Differentiating nonsmooth networks: Kakade and Lee, \emph{Provably
  Correct Automatic Subdifferentiation for Qualified Programs}, NeurIPS
  (2018); Bolte and Pauwels, \emph{Conservative Set Valued Fields,
  Automatic Differentiation, Stochastic Gradient Methods and Deep
  Learning}, Mathematical Programming (2021).
\item
  Rotary position encodings: Su et al., \emph{RoFormer: Enhanced
  Transformer with Rotary Position Embedding}, arXiv:2104.09864 (2021).
\end{itemize}
